\documentclass[12pt]{article}

% Compile with LuaLaTeX

\usepackage[margin=1.25in]{geometry}
\usepackage{setspace}
\usepackage{microtype}
\usepackage{titlesec}
\usepackage{parskip}
\usepackage{amsmath}
\usepackage{amssymb}
\usepackage{longtable}
\usepackage{array}

\newcolumntype{A}{>{\begin{otherlanguage}{arabic}\RL}p{3.2cm}<{\end{otherlanguage}}}

\usepackage{hyperref}
\usepackage{tipa}

\onehalfspacing
\usepackage{fontspec}
\usepackage[bidi=basic, main=english]{babel}
\babelprovide[import]{arabic}

\babelfont{rm}{TeX Gyre Termes}
\babelfont[arabic]{rm}[
  Renderer=HarfBuzz,
  Path=fonts/,
  Extension=.ttf,
  UprightFont=*-Regular,
  ItalicFont=*-Italic,
  BoldFont=*-Bold,
  BoldItalicFont=*-BoldItalic,
  Script=Arabic,
  Scale=1.1
]{Amiri}

\newcommand{\AR}[1]{\foreignlanguage{arabic}{#1}}

\hypersetup{
  unicode=true,
  colorlinks=true,
  linkcolor=black,
  urlcolor=blue,
  citecolor=black
}

\title{\textbf{Nigerian Pidgin English:\\
Grammar and Core Vocabulary with Script Representation}}
\author{Flyxion}
\date{\today}

\begin{document}

\maketitle

\begin{abstract}
\noindent
Nigerian Pidgin English is a major contact language spoken by tens of millions of people across Nigeria and neighboring regions of West Africa. Although frequently described in popular discourse as a “pidgin” in the sense of a simplified or auxiliary code, it exhibits systematic phonology, stable morphosyntax, productive aspect marking, and discourse strategies comparable to other creole and post-creole continua. Its grammatical organization cannot be reduced to defective English; rather, it reflects a structured outcome of sustained multilingual contact.

This essay offers a formal introduction to Nigerian Pidgin as a rule-governed linguistic system. It surveys its segmental phonology, pronominal paradigm, tense–aspect system, serial verb constructions, negation patterns, interrogative formation, and high-frequency lexical domains. In addition to descriptive analysis, the study develops parallel representational systems in both Latin phonemic transcription and a constrained Arabic-script transliteration framework. The goal is not conversational fluency but structural orientation: to demonstrate that Nigerian Pidgin possesses internal regularities and distributional constraints sufficient to sustain grammatical description, cross-script encoding, and computational modeling.
\end{abstract}

\newpage
\section{Introduction}

Languages that arise from contact are often mischaracterized as transitional or structurally incomplete. Such descriptions obscure the fact that contact languages stabilize through regularization, constraint, and functional pressure. Once stabilized, they display patterned morphosyntax and systematic phonology no less than historically older literary languages. Nigerian Pidgin is exemplary in this regard.

The purpose of this study is not to rehearse sociolinguistic debates over classification, nor to provide a phrasebook-style inventory of expressions. Instead, it adopts a structural perspective. Nigerian Pidgin is treated here as a linguistic system whose components can be described in terms of phoneme inventories, morphosyntactic categories, and combinatorial rules. Its apparent transparency to speakers of English masks underlying reorganizations of tense–aspect marking, argument structure, and discourse management that are neither accidental nor ad hoc.

The approach taken in the following sections proceeds from foundational description to grammatical architecture. Phonological inventory establishes the sound system. Pronominal forms and aspectual markers reveal core syntactic organization. Serial verb constructions illustrate clause-level compositionality. Negation and interrogative formation demonstrate rule-governed transformations rather than lexical substitution. Vocabulary is presented not as a memorization list but as evidence of semantic structuring within recurrent domains.

By foregrounding structural regularity, this introduction aims to reposition Nigerian Pidgin within linguistic analysis as a coherent grammar rather than a derivative code. Such treatment is essential not only for descriptive accuracy but also for orthographic and computational extension. Accordingly, the study develops parallel Latin phonemic representations and a formally constrained Arabic-script transliteration system. These representational layers do not precede structure; they depend upon it. The primary object of study remains the internal organization that makes systematic cross-script encoding possible.

\section{Historical and Linguistic Background}

\subsection{Origins and Stabilization}

Nigerian Pidgin emerged from sustained contact between English-speaking traders and diverse West African language communities in the Niger Delta beginning in the sixteenth and seventeenth centuries, with Portuguese maritime contact antedating British commercial dominance \cite{Holm1988}. It developed initially as a trade and interethnic communication language and gradually stabilized across regions through what contact linguists describe as feature selection under functional pressure \cite{Mufwene2001}. While English supplies the bulk of the lexical material, the grammar reflects substrate influence from Niger--Congo languages, particularly in aspect marking, serial verb constructions, pronoun systems, and conceptual organization of the lexicon.

The process of stabilization was not uniform. Early-contact varieties were confined to coastal trading settlements; inland diffusion accelerated during the colonial period as labor migration and military recruitment spread the language across ethnic boundaries. This trajectory --- from restricted trade pidgin to general-purpose lingua franca --- distinguishes Nigerian Pidgin from short-lived contact varieties that did not survive the commercial contexts that generated them. By the late twentieth century, it had acquired native speakers in urbanized communities in Warri, Port Harcourt, Sapele, and Benin City, placing it on the pidgin--creole continuum in ways that complicate strict categorical assignment \cite{Tosco2000}.

\subsection{Speaker Population and Sociolinguistic Reach}

The scale of Nigerian Pidgin use is substantial. Faraclas estimated the speaker population at over seventy million in the 1990s \cite{Faraclas1996}; more recent assessments place the total user population --- including both first-language speakers and fluent second-language users --- at approximately one hundred and sixteen million, representing roughly half of Nigeria's population \cite{LewisSimonsFennig2016}. It is, by this count, the most widely spoken language in Nigeria after English, and the fastest-spreading variety in the country. Its geographic reach extends beyond Nigeria into Cameroon, Equatorial Guinea, and diaspora communities across Europe and North America.

Today Nigerian Pidgin functions as a lingua franca in urban centers such as Lagos, Port Harcourt, Warri, and Benin City. It is present in music, radio broadcasting, informal media, comedy, political campaigning, religious settings, and everyday interaction across class lines. Dedicated broadcast media in Nigerian Pidgin include Wazobia FM and Wazobia TV, which operate exclusively in the language, as well as a West African Pidgin English service launched by the BBC. The Wycliffe Bible Society has produced a full Nigerian Pidgin translation, published under the title \textit{Naijá Baíbul}, which extends from Matthew through Revelation --- itself a demonstration of the language's capacity to sustain sustained formal register \cite{Faraclas2013}.

\subsection{Attitudes, Stigma, and Prestige}

Despite its demographic prevalence, Nigerian Pidgin has historically been subject to institutional stigma. It has been characterized in popular and educational discourse as ``broken English,'' associated with low educational attainment, and excluded from formal schooling. This characterization has substrate in colonial language ideology, which framed proximity to Standard English as a measure of social value and conversely framed deviation as deficiency. The consequences were material: Nigerian Pidgin was forbidden in many schools, discouraged in middle-class households, and denied access to official registers.

These attitudes, however, have shifted considerably since the 1990s. Research among educated Nigerians, including university students and academic staff, finds widespread positive orientation toward the language. A survey of five hundred and ten university respondents across Nigeria's six geopolitical zones found that approximately 79 percent held positive attitudes toward Nigerian Pidgin, while roughly 77 percent characterized it as a language belonging to all Nigerians rather than to any particular class or ethnic group \cite{AjaKilani2024}. The language's ethnic neutrality --- the fact that it belongs to no single ethnic constituency and therefore does not activate interethnic competition in the way that Hausa, Yoruba, or Igbo do when proposed as lingua francas --- is frequently cited as its most important sociolinguistic asset.

Prestige, however, is stratified by gender. Research in the Nigerian diaspora in Canada found a consistent pattern in which male speakers embraced Nigerian Pidgin as a marker of ethnic solidarity and heritage identity, while female speakers more frequently aligned with Standard English as a marker of educational attainment and cosmopolitan aspiration \cite{Affia2023}. This pattern mirrors Trudgill's analysis of covert and overt prestige in contact situations: men respond to the vernacular solidarity norms that Nigerian Pidgin indexes, while women respond to the overt institutional prestige that Standard English commands. Neither orientation is linguistically irrational; they reflect differential positioning within intersecting social hierarchies.

\subsection{Policy Context and Official Status}

Nigeria presents an unusual policy situation. With over five hundred documented languages, a colonial-era official language (English), three designated national languages (Hausa, Yoruba, Igbo) whose designation has itself generated political tension, and a widely spoken de facto lingua franca (Nigerian Pidgin) that lacks any formal recognition, the country has operated without a coherent language policy for most of its post-independence history \cite{OjoKilani2020}.

Proposals for elevating Nigerian Pidgin to official status alongside English have grown more prominent in academic and policy discussions. Proponents argue that official recognition would expand access to government services for populations with limited English proficiency, enhance mass literacy campaigns, reduce the symbolic violence of an official language most citizens cannot speak fluently, and provide a politically neutral common medium that does not privilege any single ethnic group \cite{Ehondor2020}. Critics raise concerns about potential harm to English-language acquisition, difficulty of implementing a standardized orthography at scale, and the risk of marginalizing smaller indigenous languages if policy attention concentrates on Pidgin.

Comparative cases from Tanzania (where Swahili functions as both national and official language), Papua New Guinea (where Tok Pisin is constitutionally recognized), and Haiti (where Haitian Creole shares official status with French) are regularly invoked to demonstrate that creolized or pidgin-derived varieties can bear the weight of formal official functions without displacing other languages or degrading their speakers' access to global communication \cite{OjoKilani2020}.

\subsection{Substrate Influence: Grammar and Conceptual Organization}

The substrate influence on Nigerian Pidgin operates at multiple levels. At the morphosyntactic level, it produces the preverbal tense--aspect--mood particle system described in Section~5, the serial verb constructions described in Section~7, and the absence of agreement morphology. These features are not random simplifications of English but positive structures that replicate patterns found in Yoruba, Igbo, Edo, Ijaw, and related Niger--Congo languages.

At the conceptual level, substrate languages shape how the Nigerian Pidgin lexicon organizes meaning. Research on Nigerian Pidgin phraseology demonstrates that multiword expressions, though composed of English-derived phonological material, follow semantic and combinatorial patterns drawn from West African conceptual schemas rather than English ones \cite{Frackiewicz2019}. The productive body-part metaphor system --- in which \textit{belle} (belly) encodes emotional and existential states, \textit{ai} (eye) encodes attention and perception, and \textit{maut} (mouth) encodes speech acts and their social consequences --- mirrors parallel structures in Hausa, Yoruba, and Igbo, not English. This divergence means that apparent lexical familiarity for English speakers conceals underlying conceptual organization that follows African rather than European patterns. A formal grammar of Nigerian Pidgin must therefore attend not only to morphosyntactic rules but to the semantic frames within which those rules operate. These structures are addressed in greater detail in the section on phraseology below.

Crucially, Nigerian Pidgin is not ``broken English.'' It is a rule-governed system. Many of its features represent simplification relative to Standard English morphology, but simplification does not imply absence of structure. Instead, complexity is reorganized, particularly in the tense--aspect system. The reorganization follows the logic of the substrate languages, not the logic of English imperfectly acquired.

\section{Phonological Tendencies}

Nigerian Pidgin pronunciation reflects both English input and substrate phonological systems. Certain consonant clusters are simplified. Word-final consonants may be reduced. Vowel quality is more stable and less diphthongized than in many English dialects.

Examples:

\begin{quote}
\textit{school} → \textit{skul} \\
\textit{friend} → \textit{fren} \\
\textit{thing} → \textit{tin}
\end{quote}

The interdental fricatives /θ/ and /ð/ are typically replaced by /t/ or /d/ respectively:

\begin{quote}
\textit{three} → \textit{tree} \\
\textit{that} → \textit{dat}
\end{quote}

Stress patterns are more even, and rhythm often reflects syllable-timed tendencies common in West African languages rather than the stress-timed rhythm typical of British English.

\section{Core Morphosyntactic Structure}

\subsection{Basic Word Order}

Nigerian Pidgin exhibits canonical Subject--Verb--Object (SVO) word order, consistent with English surface structure but lacking the inflectional morphology characteristic of Standard English.

\begin{quote}
\textit{I chop rice.} \\
(I eat rice.)
\end{quote}

Verbs do not inflect for person or number. There is no subject--verb agreement morphology:

\begin{quote}
\textit{I go.} \\
\textit{You go.} \\
\textit{Dem go.}
\end{quote}

The absence of inflection does not reduce grammatical clarity because temporal and aspectual distinctions are encoded analytically through preverbal particles.

\subsection{Pronoun System}

The pronominal system is stable and invariant across most dialectal varieties:

\begin{center}
\begin{tabular}{ll}
\textbf{Singular} & \textbf{Plural} \\
I & We \\
You & Una \\
E (he/she/it) & Dem \\
\end{tabular}
\end{center}

Notable structural features of Nigerian Pidgin further illustrate its grammatical economy and functional distinctiveness. One of the most salient characteristics is the absence of gender distinction in the third person singular. The form \textit{e} serves uniformly for referents that, in Standard English, would require differentiation between \textit{he}, \textit{she}, or \textit{it}. This neutralization of grammatical gender reflects a streamlined pronominal system in which semantic interpretation is resolved contextually rather than morphologically.

Equally significant is the second person plural form \textit{una}, which encodes number overtly in a way that contemporary Standard English no longer does. Whereas English relies on contextual inference or dialectal innovations such as \textit{you all}, Nigerian Pidgin maintains a dedicated plural pronoun, thereby preserving a number distinction that enhances pragmatic clarity in address.

The third person plural form \textit{dem} likewise demonstrates multifunctionality. In addition to serving as an independent plural pronoun, it may operate as a post-nominal plural marker, as in constructions equivalent to ``the children them.'' This dual role exemplifies a broader tendency within the language toward analytic strategies and flexible category boundaries, where lexical items may extend beyond narrowly delimited grammatical functions.

Example:

\begin{quote}
\textit{Dem dey come.} \\
(They are coming.)
\end{quote}

Plural marking may also appear post-nominally:

\begin{quote}
\textit{Di boy dem.} \\
(The boys.)
\end{quote}

Here \textit{dem} functions as a pluralizer rather than a pronoun.

\subsection{Copula and Existential Constructions}

The copular system differs from English. The verb \textit{be} does not function as an inflected auxiliary. Instead, several strategies appear.

\textbf{Zero copula:}

\begin{quote}
\textit{I happy.} \\
(I am happy.)
\end{quote}

\textbf{Na-cleft structure:}

\begin{quote}
\textit{Na doctor e be.} \\
(It is a doctor.)
\end{quote}

\textit{Na} serves as a focus or identificational marker.

\textbf{Existential construction:}

\begin{quote}
\textit{Get problem.} \\
(There is a problem.)
\end{quote}

or

\begin{quote}
\textit{E get problem.}
\end{quote}

\subsection{Verb Invariance}

Verbs in Nigerian Pidgin do not inflect for tense or agreement. Instead, temporal reference is either contextually inferred or expressed through preverbal markers, discussed in the next section.

Examples:

\begin{quote}
\textit{I chop yesterday.} \\
(I ate yesterday.)

\textit{I go tomorrow.} \\
(I will go tomorrow.)
\end{quote}

Temporal adverbs frequently carry the burden of time specification in the absence of morphological tense.

\section{Tense Aspect and Mood}

Nigerian Pidgin relies primarily on preverbal particles rather than inflectional morphology to express tense, aspect, and modality. These particles occur between the subject and the verb and remain invariant across persons.

\subsection{Progressive Aspect: \textit{dey}}

The particle \textit{dey} marks progressive or ongoing action.

\begin{quote}
\textit{I dey chop.} \\
(I am eating.)

\textit{Dem dey work.} \\
(They are working.)
\end{quote}

The progressive is thus analytic rather than inflectional. The verb itself remains unmodified.

\subsection{Past Marker: \textit{bin}}

The particle \textit{bin} marks anterior or past reference.

\begin{quote}
\textit{I bin go.} \\
(I went / I had gone.)

\textit{Dem bin see am.} \\
(They saw him/her/it.)
\end{quote}

In many contexts, past time may also be inferred from adverbials without overt marking. The use of \textit{bin} tends to emphasize temporal distance or narrative sequencing.

\subsection{Future Marker: \textit{go}}

The particle \textit{go} expresses future reference or intention.

\begin{quote}
\textit{I go come tomorrow.} \\
(I will come tomorrow.)

\textit{E go rain.} \\
(It will rain.)
\end{quote}

Unlike English “will,” \textit{go} does not inflect and does not require auxiliary inversion in questions.

\subsection{Completive Aspect: \textit{don}}

The particle \textit{don} marks completed action, roughly corresponding to the English present perfect or perfective.

\begin{quote}
\textit{I don chop.} \\
(I have eaten / I already ate.)

\textit{Dem don finish.} \\
(They have finished.)
\end{quote}

\textit{Don} emphasizes completion rather than temporal location.

\subsection{Habitual Aspect}

Habitual meaning may be expressed through context or through repeated use of \textit{dey} in certain dialects:

\begin{quote}
\textit{I dey go market every day.} \\
(I go to the market every day.)
\end{quote}

The distinction between progressive and habitual is often context-driven rather than morphologically encoded.

\subsection{Modal Constructions}

Modal meanings are typically expressed lexically rather than through auxiliary inflection.

\textbf{Ability:}

\begin{quote}
\textit{I fit do am.} \\
(I can do it.)
\end{quote}

\textbf{Obligation:}

\begin{quote}
\textit{You for go.} \\
(You should have gone / You ought to go.)
\end{quote}

\textbf{Necessity:}

\begin{quote}
\textit{You must go.} \\
(You must go.)
\end{quote}

The system is thus largely analytic. Grammatical categories are expressed through separate lexical items rather than bound morphology.

\subsection{Particle Ordering}

Particles may combine in structured sequences. For example:

\begin{quote}
\textit{I bin dey work.} \\
(I was working.)

\textit{Dem don dey wait.} \\
(They have started waiting / They are already waiting.)
\end{quote}

The ordering reflects temporal scope, with past or completive markers typically preceding progressive marking.

\section{Negation and Interrogative Structure}

\subsection{Negation}

Negation in Nigerian Pidgin is typically expressed with the preverbal particle \textit{no}.

\begin{quote}
\textit{I no go.} \\
(I did not go / I do not go.)

\textit{Dem no dey work.} \\
(They are not working.)
\end{quote}

Negation precedes aspectual markers when both are present:

\begin{quote}
\textit{I no bin see am.} \\
(I did not see him/her/it.)
\end{quote}

The structure remains analytic and invariant; the verb itself does not change form.

\subsection{Negative Imperatives}

Prohibition is typically formed using \textit{no} or \textit{make no}:

\begin{quote}
\textit{No go there.} \\
(Don't go there.)

\textit{Make you no talk.} \\
(Do not speak.)
\end{quote}

\subsection{Yes--No Questions}

Yes--no questions are generally formed through intonation rather than auxiliary inversion.

\begin{quote}
\textit{You dey come?} \\
(Are you coming?)

\textit{Dem don finish?} \\
(Have they finished?)
\end{quote}

There is no inversion comparable to English “Are you coming?” The surface structure remains declarative, and interrogative force is conveyed prosodically.

\subsection{Wh-Questions}

Interrogative words are typically fronted but do not trigger subject--auxiliary inversion.

\begin{center}
\begin{tabular}{ll}
Wetyn & What \\
Who & Who \\
Where & Where \\
When & When \\
Why & Why \\
How & How \\
\end{tabular}
\end{center}

Examples:

\begin{quote}
\textit{Wetyn you dey do?} \\
(What are you doing?)

\textit{Where dem go?} \\
(Where did they go?)
\end{quote}

Again, the verb remains uninflected.

\subsection{Focus and Emphasis with \textit{na}}

The particle \textit{na} serves as a focus marker and clefting device.

\begin{quote}
\textit{Na me do am.} \\
(It was I who did it.)

\textit{Na today we go start.} \\
(It is today that we will begin.)
\end{quote}

The structure resembles cleft constructions in other creole languages and in certain Niger--Congo substrate languages.

\section{Serial Verb Constructions}

Serial verb constructions are common and reflect substrate influence from West African languages. Multiple verbs may appear in sequence without conjunctions, expressing closely linked actions.

\begin{quote}
\textit{I go buy book.} \\
(I will go and buy a book.)

\textit{Dem carry come.} \\
(They brought it.)

\textit{I take knife cut am.} \\
(I used a knife to cut it.)
\end{quote}

In the final example, \textit{take} functions as an instrumental marker rather than a fully independent verb in the English sense.

Serial constructions allow compact expression of direction, instrumentality, causation, and sequential action without subordination or conjunction.

\section{Core Vocabulary and Lexical Domains}

Although much of Nigerian Pidgin vocabulary derives historically from English, lexical meaning often diverges in scope, usage, or pragmatic force. The lexicon also incorporates substrate and regional elements.

\subsection{Basic Verbs}

A small set of high-frequency verbs carries much of the communicative load:

\begin{center}
\begin{tabular}{ll}
\textit{chop} & eat \\
\textit{go} & go \\
\textit{come} & come \\
\textit{see} & see \\
\textit{hear} & hear \\
\textit{talk} & speak \\
\textit{carry} & carry / bring \\
\textit{take} & take / use \\
\textit{give} & give \\
\textit{make} & make / cause \\
\end{tabular}
\end{center}

Many verbs have broader semantic range than their English equivalents. For example, \textit{chop} may also mean ``consume'' metaphorically:

\begin{quote}
\textit{Dem chop money.} \\
(They embezzled money.)
\end{quote}

\subsection{Common Nouns}

\begin{center}
\begin{tabular}{ll}
\textit{pikin} & child \\
\textit{woman} & woman \\
\textit{man} & man \\
\textit{house} & house \\
\textit{money} & money \\
\textit{wahala} & trouble \\
\textit{palava} & problem / dispute \\
\end{tabular}
\end{center}

\textit{Wahala} and \textit{palava} are widely used to indicate difficulty or complication.

\subsection{Function Words and Particles}

Function words play a critical structural role.

\begin{center}
\begin{tabular}{ll}
\textit{dey} & progressive marker \\
\textit{don} & completive \\
\textit{bin} & past \\
\textit{go} & future \\
\textit{no} & negation \\
\textit{na} & focus marker \\
\textit{wey} & relative marker \\
\end{tabular}
\end{center}

Relative clauses are typically formed using \textit{wey}:

\begin{quote}
\textit{Di man wey come yesterday.} \\
(The man who came yesterday.)
\end{quote}

\subsection{Discourse Markers}

Nigerian Pidgin includes a range of discourse markers used for emphasis, alignment, and pragmatic nuance:

\begin{center}
\begin{tabular}{ll}
\textit{sha} & mild emphasis / concession \\
\textit{o} & emphasis or emotional marking \\
\textit{abeg} & please / request softener \\
\textit{now} & intensifier or focus particle \\
\end{tabular}
\end{center}

Examples:

\begin{quote}
\textit{I go come sha.} \\
(I will come, anyway.)

\textit{Abeg help me.} \\
(Please help me.)

\textit{Na today now.} \\
(It is today indeed.)
\end{quote}

These markers operate pragmatically rather than syntactically, contributing to tone and stance.

\subsection{Word-Formation Processes}

The Nigerian Pidgin lexicon is not a static inventory borrowed wholesale from English. It is generatively productive, expanding through a range of morphological and phonological processes that operate on existing material to produce new forms. These processes, documented extensively in the spoken and written record, reveal the language's structural independence from its English etymological source \cite{AmadiKilani2021}.

\paragraph{Abbreviation and Truncation.}
Frequently used words may be reduced to their initial syllables or letters. This process is especially active in registers associated with youth, university, and urban informal speech.

\begin{center}
\begin{tabular}{lll}
\textit{Full form} & \textit{Truncated} & \textit{Meaning} \\
\hline
transport & \textit{trans} & transport \\
generator & \textit{gen} & generator \\
information & \textit{info} & information \\
\end{tabular}
\end{center}

Single-letter abbreviations also function independently: \textit{H} may stand for \textit{hungry} or \textit{hundred} depending on context, relying on use-theoretic disambiguation.

\paragraph{Initialism and Acronymy.}
Multi-word phrases may be reduced to initial letters or syllables forming pronounceable tokens:

\begin{center}
\begin{tabular}{lll}
\textit{Form} & \textit{Expansion} & \textit{Meaning} \\
\hline
JJC & Johnny Just Come & newcomer, na\"{i}f \\
OYO & On Your Own & fending for oneself \\
COC & Code of Conduct & membership marker (cultist register) \\
\end{tabular}
\end{center}

\paragraph{Blending.}
Two forms may be fused into a single unit whose phonological material combines portions of both sources:

\begin{center}
\begin{tabular}{lll}
\textit{Blend} & \textit{Source} & \textit{Meaning} \\
\hline
\textit{akanchi} & Yoruba substrate & stingy person \\
\textit{blockhead} & English compound & ignorant or obstructive person \\
\end{tabular}
\end{center}

\paragraph{Reduplication.}
Full or partial repetition of a stem produces intensified, distributive, or iterative meanings. This process, discussed further under grammatical enrichment, is among the most productive morphological operations in the language:

\begin{center}
\begin{tabular}{lll}
\textit{Reduplicated form} & \textit{Meaning} & \textit{Function} \\
\hline
\textit{wella-wella} & very good, excellent & intensification \\
\textit{small-small} & little by little & iterative / distributive \\
\textit{sofri-sofri} & slowly, gently & manner modification \\
\textit{bia-bia} & beard & noun formation \\
\end{tabular}
\end{center}

\paragraph{Symbolic and Numerical Coding.}
Certain lexical items are encoded through numerical or color symbolism, often originating in specific sociolects before generalizing:

\begin{center}
\begin{tabular}{lll}
\textit{Symbol} & \textit{Referent} & \textit{Origin} \\
\hline
404 & dog meat & Peugeot 404 pickup trucks used by dog-meat sellers \\
411 & latest news or information & North American slang borrowing \\
247 & continuously, always & from ``24/7'' \\
\textit{White} / \textit{Sky} & fifty naira note & color of the note \\
\end{tabular}
\end{center}

This last category illustrates that word-formation in Nigerian Pidgin extends beyond phonological operations into semiotic substitution: the sign vehicle may be numeric or visual rather than phonemic.

\paragraph{Conversion and Recategorization.}
Words shift grammatical category without overt morphological marking. The verb \textit{baz} (to relax) may function as a noun (relaxation, downtime); \textit{bulala} (to beat) as a noun (a beating implement). This zero-derivation follows the general analytic logic of the language, in which category is determined by syntactic position rather than bound morphology.

\paragraph{Onomatopoeia and Phonetic Motivation.}
Some items in specialized registers draw on sound-symbolic or phonetic associations. The university register \textit{jack} (to study intensively) and the social register \textit{jelof} (to enjoy oneself) exploit phonetic resemblance to create mnemonically anchored meanings. Similarly, terms such as \textit{Keziah} (a woman of loose reputation) and \textit{Lion}/\textit{Linus} (a liar) exploit the phonetic overlap between the name and the concept it encodes.

\paragraph{Register-Specific Borrowing.}
Beyond general-purpose word-formation, individual sociolects develop dedicated vocabularies through internally regulated borrowing and coinage. The university register produces forms such as \textit{expo} (examination cheat material, from ``exposure''), \textit{jambite} (a first-year student, from ``JAMB,'' the national university entrance board), and \textit{jack} (to study intensely). The motor-park register includes \textit{pilot} (driver), \textit{one nyarsh} (one seat remaining), and \textit{Wazobia} (fifty naira, from the combined name of the three major ethnic radio stations). The market trader register includes \textit{jara} (a gratuitous extra quantity given to a buyer), \textit{eleme} (onlookers who do not purchase), and \textit{Bayelsa} (a serious buyer --- motivated by a perceived association with wealth in that oil-producing state). These register-bounded vocabularies are not opaque to general Nigerian Pidgin speakers at the level of grammar, but their lexical items require familiarity with the relevant social domain. They demonstrate that the language's lexical productivity operates under socially distributed rather than uniformly accessible conditions.

\section{Semantic Extension and Pragmatic Usage}

Lexical meaning in Nigerian Pidgin frequently extends beyond its etymological English base. Words may broaden, narrow, or shift in pragmatic value.

For example:

\begin{quote}
\textit{How far?} \\
(Literal: How far?) \\
Meaning: How are things? What is happening?
\end{quote}

\begin{quote}
\textit{You dey?} \\
(Are you present? / Are you available?)
\end{quote}

Meaning often depends on context, shared cultural understanding, and prosodic cues rather than morphological marking.

\section{Phraseology and Idiom Formation}

\subsection{Substrate Conceptual Schemas in the Lexicon}

The study of Nigerian Pidgin phraseology reveals a systematic tension between surface form and underlying structure. Multiword expressions are composed, for the most part, of phonological material derived from English; yet the combinatorial and conceptual patterns they instantiate are drawn from West African substrate languages rather than from English idiom \cite{Frackiewicz2019}. This gap between etymological source and conceptual organization is one of the most structurally significant features of Nigerian Pidgin as a contact language. It means that apparent transparency to English speakers is partly illusory: the words are recognizable, but the semantic frames within which they combine are not.

The theoretical basis for this analysis is use-theoretic: meaning is constituted by the conditions and contexts of use within a community, not by formal derivation from etymological source. A phrase such as \textit{belle sweet} is not interpretable by analogy with English ``belly'' and ``sweet''; its meaning (happiness, contentment, the feeling that circumstances are favorable) is determined by the conventions of use operative in Nigerian Pidgin communities \cite{Frackiewicz2019}. That those conventions mirror Hausa \textit{ciki ya daɗa} and Yoruba emotional idiom rather than English body-metaphor patterns reflects the substrate's continuing organizational role in the language's semantic architecture.

\subsection{Body-Part Metaphor Systems}

The most productive domain of Nigerian Pidgin phraseology is body-part metaphor. Specific anatomical terms function as conceptual anchors for entire fields of meaning, each term governing a family of related expressions through a consistent mapping between physical domain and abstract domain. This organization is shared with substrate languages and diverges systematically from English metaphorical conventions.

\paragraph{\textit{Belle} (Belly / Stomach).}
\textit{Belle} serves as the primary locus of emotional and existential experience in Nigerian Pidgin phraseology. The belly encodes emotional states, social relations, and life conditions in ways that have no direct English equivalent.

\begin{center}
\begin{tabular}{lll}
\textit{Expression} & \textit{Literal gloss} & \textit{Meaning} \\
\hline
\textit{belle sweet} & belly is sweet & happiness, contentment \\
\textit{belle do am} & belly does it to him/her & be pregnant \\
\textit{belle full} & belly is full & satisfied; also pregnant \\
\textit{belle no good} & belly is not good & envy, ill-wishing \\
\textit{get belle} & have belly & to be pregnant \\
\end{tabular}
\end{center}

The polysemy of \textit{belle} --- covering both the physical organ and pregnancy as a state --- is not accidental; it mirrors the Hausa and Igbo conceptual organization in which abdominal fullness and gestation are encoded by the same or closely related forms. The extension to emotional states reflects a broader West African metaphorical schema in which internal bodily sensations ground emotional and social experience.

\paragraph{\textit{Ai} (Eye).}
\textit{Ai} governs expressions related to perception, attention, social orientation, and desire. The eye encodes both directed perception and the evaluative, covetous, or protective attention one person pays to another.

\begin{center}
\begin{tabular}{lll}
\textit{Expression} & \textit{Literal gloss} & \textit{Meaning} \\
\hline
\textit{open ai} & open eye & to be alert, watchful \\
\textit{ai go} & eye goes & to desire, to covet \\
\textit{ai no gree} & eye does not agree & to be dissatisfied with what one sees \\
\textit{ai red} & eye is red & anger (the angry eye) \\
\end{tabular}
\end{center}

The use of \textit{ai} for desire (\textit{ai go}) mirrors Yoruba and Hausa constructions in which the eye's movement toward an object encodes wanting rather than mere perception. This contrasts with Standard English, where desire is typically expressed through heart metaphors or lexically.

\paragraph{\textit{Maut} (Mouth).}
\textit{Maut} governs expressions related to speech, social communication, deception, and the power of utterance. Speech acts are conceptualized through the physical configuration and movement of the mouth.

\begin{center}
\begin{tabular}{lll}
\textit{Expression} & \textit{Literal gloss} & \textit{Meaning} \\
\hline
\textit{maut no good} & mouth is not good & tendency to gossip or slander \\
\textit{maut dey} & mouth is there & to be talkative \\
\textit{open maut} & open mouth & to speak openly, to reveal \\
\textit{tie maut} & tie mouth & to keep a secret \\
\end{tabular}
\end{center}

\subsection{Verb-Based Compound Formation}

Beyond body-part idiom, Nigerian Pidgin produces multiword expressions through verb-based compounding. Several structural types are attested.

\paragraph{Verb--Verb Compounds.}
Two verbs combine to form a single semantic unit encoding a complex or resultative action:

\begin{center}
\begin{tabular}{lll}
\textit{Compound} & \textit{Component verbs} & \textit{Meaning} \\
\hline
\textit{die wake} & die + wake & recover financially after ruin \\
\textit{carry go} & carry + go & take away \\
\textit{fall down} & fall + down & collapse \\
\end{tabular}
\end{center}

The compound \textit{die wake} illustrates conceptual blending: the sequence of dying and waking maps onto a cycle of financial collapse and recovery. The meaning is not derivable from the sum of the verbs' individual meanings in English but is motivated by a cultural schema in which resurgence after apparent defeat is a significant social concept.

\paragraph{Verb--Noun Compounds.}
A verb combines with its direct object to form a lexicalized expression:

\begin{center}
\begin{tabular}{lll}
\textit{Compound} & \textit{Gloss} & \textit{Meaning} \\
\hline
\textit{chop bullet} & eat bullet & to be shot \\
\textit{chop money} & eat money & to embezzle \\
\textit{carry belle} & carry belly & to be pregnant \\
\end{tabular}
\end{center}

The productivity of \textit{chop} as a verb of consumption extending to non-alimentary domains (violence, financial appropriation) reflects its broader semantic scope relative to English ``eat.''

\paragraph{Noun--Noun Compounds.}
Two nouns combine to form a classificatory or relational expression:

\begin{center}
\begin{tabular}{lll}
\textit{Compound} & \textit{Gloss} & \textit{Meaning} \\
\hline
\textit{man pikin} & man child & male child, son \\
\textit{woman pikin} & woman child & female child, daughter \\
\textit{papa battalion} & father battalion & man with many children \\
\end{tabular}
\end{center}

The compound \textit{papa battalion} is particularly expressive: the military register of \textit{battalion} is recruited to describe domestic proliferation, producing a humorous and emphatic intensification of \textit{papa} (father). This type of semantic surprise is a characteristic feature of Nigerian Pidgin idiomatic creativity.

\subsection{Substrate Calques}

Some Nigerian Pidgin expressions are structural calques from specific substrate languages: the form is composed of Nigerian Pidgin words, but the combinatorial structure replicates a specific construction from Hausa, Yoruba, or another substrate language.

\begin{center}
\begin{tabular}{lll}
\textit{NPE expression} & \textit{Substrate source} & \textit{Meaning} \\
\hline
\textit{chop cockroach} & Hausa \textit{ci wākē} (eat bean/pea) & to be pregnant (euphemistic) \\
\textit{put mind to body} & Yoruba \textit{kó iyè sí ara} & to pay attention, focus \\
\end{tabular}
\end{center}

These calques demonstrate that the substrate's influence on Nigerian Pidgin phraseology is not limited to phonological or morphosyntactic transfer but extends to specific idiomatic constructions that have been relexified into the English-derived phonological material of the pidgin.

\subsection{Polysemy and Contextual Disambiguation}

Nigerian Pidgin lexical items frequently carry multiple, sometimes distant, senses that are disambiguated by context. The verb \textit{die}, for instance, extends beyond literal death to encode intensification (\textit{I die laugh} = I laughed intensely), resistance (\textit{e no go die} = it will not give way), and completion. The word \textit{belle} covers both the anatomical stomach and the social fact of pregnancy. The word \textit{chop} covers eating, sexual consumption, financial misappropriation, and being struck by a projectile.

This systematic polysemy is not lexical ambiguity in the English sense of accidentally homophonous forms. It is motivated extension within coherent semantic domains, governed by consistent conceptual metaphors. The grammar of disambiguation is contextual and pragmatic rather than morphological: no inflection marks which sense of \textit{chop} is active. This places a high burden on shared context and cultural knowledge, and it explains why apparent intelligibility of Nigerian Pidgin for English speakers is frequently illusory at the level of idiomatic expression even when it holds at the level of grammatical structure.

\section{Grammatical Enrichment}

\subsection{Reduplication}

Reduplication functions as an intensifier, distributive marker, or expressive device. It is common in many West African languages and appears productively in Nigerian Pidgin.

\textbf{Intensification:}

\begin{quote}
\textit{small small} \\
(little by little)

\textit{fine fine} \\
(very fine / extremely beautiful)
\end{quote}

\textbf{Distributive meaning:}

\begin{quote}
\textit{Dem come one one.} \\
(They came one by one.)
\end{quote}

Reduplication adds semantic nuance without inflectional morphology.

\subsection{Comparative and Superlative Structures}

Comparison does not rely on inflectional endings such as “-er” or “-est.” Instead, comparison is expressed analytically.

\textbf{Comparative with \textit{pass}:}

\begin{quote}
\textit{Dis one big pass dat one.} \\
(This one is bigger than that one.)
\end{quote}

\textit{Pass} functions as a comparative marker meaning “more than.”

\textbf{Superlative with \textit{pass all}:}

\begin{quote}
\textit{Na im fine pass all.} \\
(He/she is the most beautiful.)
\end{quote}

Comparison is thus relational rather than morphologically marked.

\subsection{Quantifiers and Plurality}

Plural marking is optional and often contextually determined. The plural marker \textit{dem} may follow a noun.

\begin{quote}
\textit{Book dem dey table.} \\
(The books are on the table.)
\end{quote}

Quantifiers include:

\begin{center}
\begin{tabular}{ll}
\textit{plenty} & many / much \\
\textit{small} & little \\
\textit{all} & all \\
\textit{every} & every \\
\end{tabular}
\end{center}

Examples:

\begin{quote}
\textit{Plenty people come.} \\
(Many people came.)

\textit{Small water remain.} \\
(A little water remains.)
\end{quote}

\subsection{Possession}

Possession is typically expressed through juxtaposition rather than apostrophe marking.

\begin{quote}
\textit{Di man house.} \\
(The man's house.)

\textit{My papa car.} \\
(My father's car.)
\end{quote}

The possessive relation is interpreted structurally rather than morphologically.

\subsection{Relative Clauses}

The relative marker \textit{wey} introduces subordinate descriptive clauses.

\begin{quote}
\textit{Di woman wey dey talk.} \\
(The woman who is speaking.)
\end{quote}

No relative pronoun inflection occurs. The marker remains invariant.

\subsection{Conditional Constructions}

Conditional statements are typically formed with \textit{if} or with serial structures.

\begin{quote}
\textit{If you come, I go happy.} \\
(If you come, I will be happy.)
\end{quote}

Hypothetical or counterfactual nuance may be conveyed through context or through use of \textit{for}.

\begin{quote}
\textit{You for tell me.} \\
(You should have told me.)
\end{quote}

\subsection{Imperatives and Causatives}

Imperatives are generally bare verbs:

\begin{quote}
\textit{Come here.}

\textit{Carry am.}
\end{quote}

Causation frequently uses \textit{make}:

\begin{quote}
\textit{Make I go.} \\
(Let me go.)

\textit{Make dem start.} \\
(Let them begin.)
\end{quote}

Here \textit{make} functions as a causative or permissive marker.

\section{Pragmatics and Politeness Strategies}

Nigerian Pidgin operates within a rich pragmatic ecology in which politeness, hierarchy, solidarity, and emotional stance are frequently conveyed through particles, repetition, tone, and lexical choice rather than through formal grammatical marking.

\subsection{Softening Devices}

The particle \textit{abeg} functions as a politeness marker, mitigating requests and commands.

\begin{quote}
\textit{Abeg give me water.} \\
(Please give me water.)
\end{quote}

It can also appear mid-sentence to soften disagreement:

\begin{quote}
\textit{Abeg, no vex.} \\
(Please, do not be angry.)
\end{quote}

\subsection{Emphatic and Affective Particles}

The sentence-final particle \textit{o} marks emotional emphasis, insistence, or affective engagement.

\begin{quote}
\textit{Come here o.} \\
(Come here, I insist / urgently.)
\end{quote}

The particle \textit{sha} may signal concession or mild contrast:

\begin{quote}
\textit{I go try sha.} \\
(I will try anyway.)
\end{quote}

\subsection{Solidarity and Alignment}

Greetings and discourse openings often establish social alignment.

\begin{quote}
\textit{How far?} \\
(How are things?)

\textit{You dey?} \\
(Are you there? / Are you around?)
\end{quote}

Such expressions function less as literal questions and more as relational markers.

\subsection{Respect and Hierarchy}

Honorific marking is generally lexical rather than morphological. Terms such as \textit{oga} (boss) or \textit{madam} index respect.

\begin{quote}
\textit{Oga, I don finish.} \\
(Boss, I have finished.)
\end{quote}

Politeness strategies rely heavily on intonation, context, and relational expectations.

\section{Code-Switching and English Overlap}

Nigerian Pidgin exists in a dynamic continuum with Standard Nigerian English and with regional languages such as Yoruba, Igbo, and Hausa. Speakers frequently alternate between varieties depending on context, audience, and social positioning.

\subsection{Lexical Mixing}

Code-switching may involve insertion of English technical vocabulary into Pidgin structure.

\begin{quote}
\textit{I go submit di assignment tomorrow.}
\end{quote}

The syntactic frame remains Pidgin while lexical items derive from Standard English.

\subsection{Register Shifting}

In formal contexts, speakers may shift toward Standard English morphology while retaining Pidgin rhythm or particles. In informal contexts, greater reliance on Pidgin structures appears.

This fluidity complicates attempts at strict grammatical boundary drawing. Nigerian Pidgin functions not as an isolated code but as part of a sociolinguistic spectrum.

\section{Regional Variation}

Although broadly intelligible across Nigeria, Nigerian Pidgin exhibits regional phonological and lexical variation.

In southern urban centers such as Warri, Port Harcourt, and Lagos, Pidgin often functions as a primary lingua franca. In northern regions, usage may be more restricted and influenced by Hausa structures.

Lexical items such as \textit{wahala} may carry slightly different pragmatic weight across regions. Pronunciation and intonation patterns vary significantly.

The language is thus internally heterogeneous rather than uniform.

\section{Orthography and Standardization}

\subsection{The Current Orthographic Situation}

Nigerian Pidgin lacks a fully standardized orthography, and this absence has consequences for both linguistic description and language planning. Spellings vary according to phonetic intuition, English convention, and individual or institutional preference, producing inconsistency across texts, broadcasts, and educational materials.

For example, the progressive marker may appear as:

\begin{quote}
\textit{dey}, \textit{de}, or \textit{day}
\end{quote}

Similarly, \textit{wetin} may appear as \textit{wetyn}, and the focus marker \textit{na} may appear as \textit{na}, \textit{nah}, or simply be omitted in informal writing. Such variation is not random: each spelling reflects a different implicit theory of the relationship between Nigerian Pidgin and English --- whether the language is being written as a modified English or as an autonomous phonemic system. The choice of spelling convention thus encodes an implicit ideological stance on the language's status.

Orthographic fluidity reflects the primarily oral historical transmission of the language and the long exclusion of Nigerian Pidgin from formal educational and institutional settings. Because the language was suppressed rather than cultivated in schools, no centralized orthographic authority had occasion to standardize it. The result is a situation analogous to the pre-standardization stages of many now-prestigious European languages: vitality without graphemic stability.

\subsection{Standardization Efforts and Institutional Context}

Despite the absence of an official orthography, standardization efforts are underway and have produced some convergence in practice. The body tasked with overseeing standardization is the Naijá Lángwèj Akadémì (NLA), an institution explicitly modeled on national language academies in other multilingual states. The NLA has undertaken vocabulary codification and orthographic recommendation work, though its authority remains limited by the absence of formal state endorsement \cite{Ehondor2020}.

In parallel, broadcast media have begun to exercise a de facto standardizing influence. Wazobia FM and Wazobia TV, which operate exclusively in Nigerian Pidgin, have necessarily adopted consistent spelling conventions for on-screen text, subtitles, and written promotional material. The BBC Pidgin service, launched for West African audiences, similarly operates under editorial style guides that impose consistency across output. These informal media standards constitute a practical orthographic baseline even without formal legislative backing.

The most significant completed standardization project is the Wycliffe Bible Society's full Nigerian Pidgin translation, known as the \textit{Naijá Baíbul}. This translation represents a sustained effort to render sustained formal and literary register in the language, and to do so in a consistent, documented orthography. A passage from the Gospel of Matthew illustrates both the capabilities and the conventions adopted:

\begin{quote}
\textit{Som Farisi pipol and law tishas kom mit Jisos for Jerusalem ask am, `Why yor disaipols nor dey obey awa eldas tradishon? Dem no dey wash dia hand bifor dem chop.'} \\
(Matthew 15:1--2)
\end{quote}

The rendering of ``Scribes and Pharisees'' as \textit{Farisi pipol and law tishas} is notable: it translates the referential content into accessible Pidgin categories rather than transliterating the proper names unchanged. This represents a translational philosophy that prioritizes communicative transparency over terminological fidelity --- the same philosophy that historically drove the production of vernacular Bibles in European languages.

\subsection{Standardization as a Policy and Linguistic Problem}

Standardization of Nigerian Pidgin faces a set of interrelated technical and sociolinguistic problems that any future orthographic authority will need to resolve.

The first problem is the choice between phonemic and etymological representation. An etymological orthography --- one that preserves the English spelling of source words --- would write \textit{know} as \textit{know} even though the /k/ is never pronounced in Nigerian Pidgin. A phonemic orthography would write \textit{no} for both the negation particle and the cognate of English ``know,'' producing homographic ambiguity but graphophonemic transparency. The analysis in this study adopts the phonemic approach, which is also the direction favored by the NLA and by most academic descriptions of the language.

The second problem is dialect variation. Nigerian Pidgin exhibits regional phonological variation --- particularly in vowel quality and the realization of certain consonant clusters --- that a single orthography cannot simultaneously represent with perfect fidelity for all speakers. The standard solution adopted in comparable cases (Swahili, Tok Pisin) is to standardize around a prestige or educated variety while acknowledging that regional pronunciation will diverge from the written norm. A seven-vowel system with /ɛ/ and /ɔ/ distinct from /e/ and /o/ would be more faithful to many speakers' phonology but would add graphemic complexity; the five-vowel baseline adopted here trades fine-grained accuracy for stability and accessibility.

The third problem is register and lexical scope. A standard orthography must decide how to handle the large registers of domain-specific vocabulary that operate in sociolects (motor-park, market-trader, university, and so on) without either excluding them from the standard or legitimizing all variants uncritically. The NLA's approach is to prioritize core grammar and high-frequency vocabulary while treating register-specific forms as a research and documentation task rather than an immediate standardization priority.

Survey data suggest that the educated Nigerian population views standardization as achievable: approximately 53 percent of five hundred and ten respondents across six geopolitical zones agreed that orthographic standardization is feasible, while 66 percent agreed that special programs and policies are needed to support the language's formal development \cite{AjaKilani2024}. These figures indicate that the obstacle to standardization is political will and institutional capacity rather than popular resistance.

\subsection{International and Comparative Precedents}

The standardization trajectory of other contact varieties is instructive. Tok Pisin, the English-based creole of Papua New Guinea, achieved constitutional recognition and orthographic standardization after a period of comparable sociolinguistic informality. Its standardization involved exactly the combination of media convention, religious translation, educational adoption, and legislative recognition that proponents of Nigerian Pidgin standardization advocate \cite{Holm1988}. Haitian Creole, standardized in 1979 and granted co-official status alongside French, demonstrates that a creolized variety emerging from colonial contact can sustain formal domains including law, education, and government without displacing associated languages.

The Ajami tradition --- the use of Arabic script for writing sub-Saharan African languages, with a history stretching across Hausa, Wolof, Mandinka, Fula, Yoruba, and others --- offers a different kind of precedent \cite{Ngom2010}. Ajami orthographies were produced by and for Muslim communities across West Africa, often without formal institutional backing, and demonstrate that alternative script systems can achieve significant stability and textual productivity in the absence of official endorsement. The Arabic-script transliteration system developed in subsequent sections of this study draws on this tradition while adapting it to the phonological requirements of Nigerian Pidgin.

\subsection{The Relationship Between Orthographic Choice and Structural Perception}

Orthography is not a neutral recording device. The script and spelling conventions adopted for a language shape how speakers, learners, and analysts perceive its structure. An orthography that preserves English spellings foregrounds Nigerian Pidgin's historical derivation from English and implicitly frames it as a modified or simplified English. An orthography built on phonemic principles foregrounds its autonomy as a system with its own sound-symbol correspondences. An Arabic-script representation severs the visual connection to English entirely and makes the language's particle-based analytic grammar directly legible at the graphemic level --- because particles appear as discrete visual units rather than being embedded in the run-on conventions of English spelling.

This is not merely a theoretical point. Computational modeling, cross-linguistic comparison, and language pedagogy are all affected by orthographic choice. A phonemically consistent orthography makes it possible to build deterministic transliteration systems, supports computational parsing, and reduces the learning burden for speakers of substrate languages who have not been educated in English orthographic conventions. The formal transliteration framework developed in subsequent sections of this study is predicated on a phonemically consistent Latin baseline, from which Arabic-script encoding follows algorithmically. The choice of baseline orthography is therefore not preliminary to the formal analysis; it is part of it.

\section{Annotated Sample Dialogue}

\begin{quote}
A: \textit{How far? You dey?} \\
B: \textit{I dey. Wetin happen?} \\
A: \textit{I bin call you yesterday but you no pick.} \\
B: \textit{Ah, I bin dey work. I don finish now.} \\
A: \textit{Good. Make we go market small small.} \\
B: \textit{Okay. I go come now.}
\end{quote}

\textbf{Analysis:}

\textit{How far?} functions as greeting rather than literal distance inquiry. \\
\textit{I dey} marks presence or availability. \\
\textit{bin} marks past action. \\
\textit{no} marks negation. \\
\textit{don} marks completion. \\
\textit{make we} forms a hortative construction. \\
\textit{go} marks future intention.

\section{Phonology and an Alternative Orthographic Proposal}

\subsection{Phonological Overview}

Nigerian Pidgin phonology reflects substrate influence from West African languages and superstrate influence from English. The sound system is generally simpler in syllable structure than Standard English, with reduced consonant clusters and strong preference for open syllables in many regional varieties.

Consonant clusters at word-final position are frequently simplified:

\begin{quote}
\textit{ask} $\rightarrow$ \textit{aks} \\
\textit{child} $\rightarrow$ \textit{chail}
\end{quote}

Vowel quality tends toward a reduced inventory compared to English, and diphthongs are often monophthongized depending on speaker background.

Tone is not lexically contrastive in the way it is in Yoruba or Igbo, but intonation carries significant pragmatic meaning, particularly in interrogatives and emphasis.

\subsection{Syllable Structure}

The dominant syllable structure approximates (C)V(C), though final consonants may be weakened or deleted in casual speech. Many borrowings from English undergo phonotactic adaptation to align with West African patterns.

Reduplication and particle-based grammar further reinforce rhythmic and prosodic regularity.

\subsection{The Case for Orthographic Reconsideration}

The current Latin-based orthography inherits English spelling conventions, which introduce inconsistency between sound and symbol. Variants such as \textit{dey}, \textit{de}, and \textit{day} illustrate instability.

An orthography more closely aligned to phonetic realization could reduce ambiguity. One possible direction is transliteration into Arabic script (Ajami-style adaptation), which has historical precedent in West Africa across Hausa and other languages.

Such a system would not imply cultural replacement but phonetic alignment. Arabic script represents consonantal structure efficiently and can be modified with diacritics to capture vowel distinctions.

\subsection{Proposed Arabic-Based Transliteration Framework}

The following is a preliminary mapping based on approximate phonetic correspondence.

Consonants:

\begin{center}
\begin{tabular}{ll}
b & \AR{ب} \\
d & \AR{د} \\
f & \AR{ف} \\
g & \AR{غ} or \AR{گ} (modified) \\
k & \AR{ك} \\
l & \AR{ل} \\
m & \AR{م} \\
n & \AR{ن} \\
p & \AR{پ} (Persian extension) \\
r & \AR{ر} \\
s & \AR{س} \\
t & \AR{ت} \\
w & \AR{و} \\
y & \AR{ي} \\
\end{tabular}
\end{center}

Vowels would be represented through short vowel diacritics and long vowel letters:

\begin{center}
\begin{tabular}{ll}
a & \AR{َ} \\
i & \AR{ِ} \\
u & \AR{ُ} \\
long a & \AR{ا} \\
long i & \AR{ي} \\
long u & \AR{و} \\
\end{tabular}
\end{center}

\subsection{Illustrative Examples}

\textit{I dey} could be rendered:

\begin{quote}
\AR{ا دي}
\end{quote}

\textit{Wetin happen?}

\begin{quote}
\AR{وَتِن هَپَن}
\end{quote}

\textit{I no go}

\begin{quote}
\AR{ا نو غو}
\end{quote}

This transliteration prioritizes phonetic transparency over etymological fidelity.

\subsection{Advantages and Challenges}

The proposed orthographic framework offers several theoretical and practical advantages. Most notably, it enables a closer alignment between phonological structure and graphic representation, thereby reducing the distance between sound and symbol that often complicates the transcription of Nigerian Pidgin. In addition, by situating the system within the broader intellectual and historical continuum of Ajami literacy traditions across West Africa, the proposal opens a pathway for continuity rather than rupture. Such alignment may facilitate cross-linguistic familiarity among communities already accustomed to Arabic-derived scripts, while also supporting the development of locally grounded literacy practices.

At the same time, significant challenges remain. Nigerian Pidgin exhibits substantial dialectal variation, particularly in its vowel realizations, and any fixed orthographic mapping risks privileging one regional norm over others. The absence of a standardized tonal marking system further complicates matters, especially in contexts where prosodic distinctions may carry semantic weight. Beyond linguistic considerations, script reform inevitably intersects with sociopolitical sensitivities. Orthography is rarely neutral; it encodes histories of power, education, religion, and identity. Any attempt to formalize or modify script usage must therefore proceed with attentiveness to community reception and cultural meaning. Finally, technical questions of encoding, font support, keyboard layouts, and digital interoperability pose nontrivial barriers to widespread adoption.

For these reasons, the present proposal is best understood as exploratory rather than prescriptive. Its aim is not cultural substitution, nor the displacement of existing Latin-based conventions, but the pursuit of phonological coherence and structural clarity within an expanded graphic repertoire.

\subsection{Theoretical Implication}

Orthography shapes perception of linguistic structure. A phonetic Arabic-based script would foreground syllabic rhythm and particle-based grammar rather than English lexical ancestry. It would visually decouple Nigerian Pidgin from Standard English, potentially reinforcing its autonomy as a linguistic system.

\section{Toward a Phonetic Ajami Orthography for Nigerian Pidgin}

\subsection{Design Principles}

The proposed Arabic-based transliteration system is not etymological but phonetic. Its primary objective is one-to-one mapping between sound and symbol, thereby minimizing the orthographic opacity inherited from English.

Three structural principles guide the system:

First, consonants are mapped according to phonetic similarity rather than English spelling convention.

Second, vowels are represented consistently through diacritics and long-vowel letters, prioritizing phonemic transparency.

Third, grammatical particles are preserved visually as independent units, reinforcing the analytic structure of Nigerian Pidgin.

\subsection{Consonant Mapping}

The consonant inventory largely overlaps with Arabic but requires extended letters for sounds not native to Classical Arabic. Persian extensions such as \AR{پ} for /p/ are employed where necessary to preserve phonemic contrast.

For example:

\begin{center}
\begin{tabular}{ll}
b & \AR{ب} \\
d & \AR{د} \\
f & \AR{ف} \\
g & \AR{غ} / \AR{گ} \\
k & \AR{ك} \\
l & \AR{ل} \\
m & \AR{م} \\
n & \AR{ن} \\
p & \AR{پ} \\
r & \AR{ر} \\
s & \AR{س} \\
t & \AR{ت} \\
w & \AR{و} \\
y & \AR{ي} \\
h & \AR{ه} \\
\end{tabular}
\end{center}

Affricates such as /ch/ and /j/ may be represented using adapted forms, for example \AR{تش} and \AR{ج}, respectively.

\subsection{Vowel Representation}

Short vowels are indicated through diacritics:

\[
a = \AR{َ} \qquad
i = \AR{ِ} \qquad
u = \AR{ُ}
\]

Long vowels are represented through \AR{ا}, \AR{ي}, and \AR{و}.

For example:

\begin{quote}
\AR{آيْ} \AR{دي} \\
(I dey)
\end{quote}

Here \AR{آيْ} captures the diphthongal realization of “I,” while \AR{دي} captures the progressive particle /deɪ/.

\subsection{Grammatical Particles in Script}

A significant structural advantage of this orthography is the visual stabilization of analytic particles.

\begin{quote}
\AR{نَا سُوْ} \\
(Na so)

\AR{دِمْ دُونْ كَمْ} \\
(Dem don come)
\end{quote}

Particles such as \AR{نَا} (na), \AR{دُونْ} (don), and \AR{دِي} (dey) become visually discrete units, reinforcing their grammatical independence.

This representation visually encodes the analytic structure of tense and aspect marking more transparently than English-based spelling.

\subsection{Phonotactic Adaptation}

The orthographic system resolves cluster reduction by reflecting actual spoken realization rather than English orthography.

For example:

\begin{quote}
\AR{وِيتِنْ} \\
(Wetin)
\end{quote}

This avoids misleading English silent letters and aligns orthography more closely with speech rhythm.

\subsection{Semantic Autonomy Through Script}

By removing English orthographic anchoring, the script visually decouples Nigerian Pidgin from Standard English. This has structural consequences.

The language appears as a system in its own right rather than a deviation from English. Particles gain structural salience. Prosodic rhythm becomes more apparent.

Orthographic reform therefore does not merely change script; it reorganizes perception of grammatical structure.

\subsection{Limitations and Technical Questions}

Dialectal vowel variation presents challenges. The distinction between /e/ and /ɛ/ or between /o/ and /ɔ/ may require additional diacritic marking if phonemic contrast is deemed significant.

Digital encoding requires standardization of extended letters such as \AR{گ} and \AR{پ} across platforms.

Sociolinguistic reception would likely vary across regions and communities.

\subsection{Algorithmic Transliteration}

The system is suitable for deterministic transliteration because it operates primarily on phoneme-to-grapheme mapping rather than etymological mapping.

A bidirectional algorithm would require:

A phonemic tokenizer for Latin Pidgin input.

A rule-based consonant mapping.

A vowel diacritic insertion rule sensitive to syllable structure.

Optional tone marking if future refinement demands it.

The resulting system is structurally simpler than English orthography because it does not encode historical spelling residue.

\section{Phoneme Inventory and IPA Mapping}

This section proposes a working phoneme inventory for Nigerian Pidgin suitable for orthographic design and then specifies an IPA-to-script mapping consistent with the transliteration practice illustrated in the corpus above. The goal is not to legislate a single ``true'' pronunciation for all regions, but to provide a coherent baseline inventory that supports deterministic transliteration.

\subsection{Working Phoneme Inventory}

Nigerian Pidgin exhibits regional variation, but a practical orthography can be built around a core consonant system and a modest vowel system. The inventory below is intended as an operational compromise: fine-grained enough to support stable transcription, coarse-grained enough to remain dialect-robust.

\paragraph{Consonants}

\[
/p,\; b,\; t,\; d,\; k,\; g,\;
m,\; n,\; \eta,\;
f,\; v,\; s,\; z,\; h,\;
w,\; j,\;
l,\; r,\;
tʃ,\; dʒ/
\]

Interdental fricatives /\texttheta/ and /\dh/ are excluded, since they are not phonemic in most Nigerian Pidgin varieties and are typically realized as alveolar stops, as reflected in forms such as \textit{dat} and \textit{tree}.

\paragraph{Vowels}

A five-vowel system provides stable coverage:

\[
/i,\; e,\; a,\; o,\; u/
\]

A seven-vowel refinement with /ɛ/ and /ɔ/ may be optionally encoded, but orthographic stability does not require it. Diphthongs are treated as vowel sequences, with /ai/ especially frequent.

\subsubsection{Latin Baseline (Phonemic)}

Before transliteration into Arabic script, phonemic forms are normalized into a narrow Latin baseline that avoids English spelling conventions. This baseline stabilizes input and prevents interference from English orthography.

\begin{center}
\begin{tabular}{lll}
IPA & Latin & Example \\
\hline
/p/ & p & pikin \\
/b/ & b & boy \\
/t/ & t & tin \\
/d/ & d & dey \\
/k/ & k & ko \\
/g/ & g & go \\
/tʃ/ & ch & chop \\
/dʒ/ & j & enjoy \\
/f/ & f & fit \\
/v/ & v & vex \\
/s/ & s & so \\
/z/ & z & zaa (dialectal) \\
/h/ & h & how \\
/m/ & m & make \\
/n/ & n & no \\
/ŋ/ & ng & sing \\
/l/ & l & lie \\
/r/ & r & far \\
/w/ & w & wetin \\
/j/ & y & you \\
\end{tabular}
\end{center}

This level is purely phonemic. English spellings such as \textit{know}, \textit{three}, or \textit{that} are excluded in favor of \textit{no}, \textit{tri}, and \textit{dat}.

\subsubsection{Arabic Script Mapping}

Arabic transliteration is derived mechanically from the Latin baseline. Only standard Arabic letters are used.

\paragraph{Consonants}

\begin{center}
\begin{tabular}{lll}
IPA & Script & Note \\
\hline
/p/ & \AR{ب} & approximated as /b/ \\
/b/ & \AR{ب} & direct \\
/t/ & \AR{ت} & direct \\
/d/ & \AR{د} & direct \\
/k/ & \AR{ك} & direct \\
/g/ & \AR{ج} & phonetic approximation \\
/f/ & \AR{ف} & direct \\
/v/ & \AR{ف} & merged with /f/ \\
/s/ & \AR{س} & direct \\
/z/ & \AR{ز} & direct \\
/h/ & \AR{ه} & direct \\
/m/ & \AR{م} & direct \\
/n/ & \AR{ن} & direct \\
/ŋ/ & \AR{نْج} & cluster representation \\
/l/ & \AR{ل} & direct \\
/r/ & \AR{ر} & direct \\
/w/ & \AR{و} & direct \\
/j/ & \AR{ي} & direct \\
/tʃ/ & \AR{تش} & digraph \\
/dʒ/ & \AR{ج} & direct \\
\end{tabular}
\end{center}

No additional Persian letters are introduced. All approximations are handled internally through clustering or phonemic merger.

\paragraph{Vowels}

Short vowels are represented by diacritics when phonological clarity requires explicit marking. The vowel /a/ is written with \AR{َ} (fatha), /i/ and /e/ with \AR{ِ} (kasra), and /u/ and /o/ with \AR{ُ} (damma).

In unstressed environments or highly predictable contexts, short vowels may be omitted in fluent writing; however, full marking ensures deterministic back-mapping to the Latin baseline and eliminates ambiguity in pedagogical contexts.

Long vowels are represented by the corresponding matres lectionis. The vowel /aː/ is written with \AR{ا}, /iː/ with \AR{ي}, and /uː/ with \AR{و}. Where contrastive vowel length is not phonemic, these long-vowel letters function primarily as quality markers rather than as duration markers.

Diphthongs are written as ordered sequences of vowel letters. The diphthong /ai/ is represented as \AR{آي}, while /au/ is written as \AR{آو}. In environments where a glottal onset is phonetically salient, initial \AR{آ} provides a stable orthographic anchor.

\subsubsection{Worked Derivations}

The full transliteration pipeline proceeds from surface form to phonological representation, from phonology to Latin baseline, and finally from the Latin baseline to Arabic script.

\paragraph{Progressive}

The form \textit{I dey} is analyzed phonologically as /ai de/, normalized to the Latin baseline \textit{ai de}, and rendered orthographically as \AR{آيْ دِي}.

\paragraph{Negation}

The expression \textit{I no know} corresponds phonologically to /ai no no/, normalizes to \textit{ai no no}, and is written as \AR{آيْ نُو نُو}.

\paragraph{Wh-question}

The form \textit{Wetin dey happen} is analyzed as /wetin de apen/, normalized to \textit{wetin de apen}, and rendered as \AR{وِيتِنْ دِي أَبِن}.

Consonant clusters are either resolved through vowel insertion consistent with attested pronunciation or indicated with suk\={u}n, as in \AR{ْ}, where pedagogically useful to preserve segmental transparency.

\subsubsection{Functional Stability}

High-frequency grammatical particles maintain fixed orthographic forms. The progressive marker \textit{dey} is written as \AR{دِي}, the completive marker \textit{don} as \AR{دُونْ}, the copular or focus particle \textit{na} as \AR{نَا}, and the negator \textit{no} as \AR{نُو}.

Stability at this level reinforces morphosyntactic transparency and ensures that analytic tense and aspect marking remains visually distinct.

\subsubsection{Orthographic Principle}

The system is strictly phonological. Each grapheme corresponds to a surface phoneme in the working inventory. English spelling conventions are excluded from derivation. Dialectal variation can be layered through optional diacritics without altering the core mapping.

The result is a closed, internally consistent transliteration framework: inventory → Latin phonemic normalization → Arabic graphemic realization. Each stage precedes the next.

\subsection{Deterministic Transliteration Rules}

This subsection specifies a deterministic procedure for transliterating a restricted Latin baseline into an Arabic-only orthography. The procedure is designed to be implementable as a finite sequence of rewrite rules. It assumes that the Latin input is already normalized to a phonetic baseline rather than English spelling.

\subsubsection{Input Normalization}

Let the input string be a sequence of tokens separated by spaces. Each token is assumed to be lower-case Latin characters drawn from the set
\[
\{a,e,i,o,u,b,d,f,g,h,j,k,l,m,n,r,s,t,w,y\}
\]
together with the digraphs \texttt{sh} and \texttt{ch} and apostrophe-free orthography. If the source uses English-residue spellings, these must be normalized first.

A minimal normalization map that removes common English residue is:

\[
\texttt{th} \mapsto \texttt{t},\quad
\texttt{ph} \mapsto \texttt{f},\quad
\texttt{ck} \mapsto \texttt{k},\quad
\texttt{qu} \mapsto \texttt{kw}.
\]

This stage is not a model of pronunciation; it is a structural stabilizer for transliteration.

\subsubsection{Tokenization}

Let the token sequence be $T_1,\dots,T_N$. Transliteration applies tokenwise. Punctuation may be retained and copied directly after the corresponding token.

\subsubsection{Consonant Mapping}

Within each token, consonantal segments are mapped deterministically by the following substitution table.

\begin{center}
\begin{tabular}{lll}
Latin & Arabic & Comment \\
\hline
b & \AR{ب} & \\
d & \AR{د} & \\
f & \AR{ف} & \\
g & \AR{ج} & Arabic-only approximation \\
h & \AR{ه} & \\
j & \AR{ج} & shares glyph with g; resolved by context \\
k & \AR{ك} & \\
l & \AR{ل} & \\
m & \AR{م} & \\
n & \AR{ن} & \\
r & \AR{ر} & \\
s & \AR{س} & \\
t & \AR{ت} & \\
w & \AR{و} & consonantal /w/ \\
y & \AR{ي} & consonantal /y/ \\
sh & \AR{ش} & \\
ch & \AR{تش} & Arabic-only cluster encoding \\
\end{tabular}
\end{center}

The map is applied greedily left-to-right with longest-match priority, so \texttt{sh} and \texttt{ch} are rewritten before single-letter consonants.

\subsubsection{Vowel Strategy}

Two regimes are possible. The first is a readability-first regime that writes vowels as letters where possible. The second is a precision-first regime that uses diacritics systematically. The transliteration system in this essay adopts the readability-first regime as default, with diacritics treated as optional refinement.

The readability-first vowel map is:

\begin{center}
\begin{tabular}{lll}
Latin & Arabic & Comment \\
\hline
a & \AR{ا} & default long \\
i & \AR{ي} & default long \\
u & \AR{و} & default long \\
e & \AR{ي} & nearest-grapheme approximation \\
o & \AR{و} & nearest-grapheme approximation \\
\end{tabular}
\end{center}

This produces a stable orthography without requiring diacritic entry. It intentionally sacrifices fine vowel quality distinctions, which in Nigerian Pidgin are often predictable from lexical context.

An optional refinement is to allow short-vowel diacritics in pedagogical mode:

\[
a \mapsto \AR{َ},\qquad i/e \mapsto \AR{ِ},\qquad u/o \mapsto \AR{ُ}
\]
when the writer wishes to avoid ambiguity.

\subsubsection{Special Lexical Tokens}

Certain high-frequency function words are best treated lexically to preserve recognizable forms and avoid overgeneration. Define a small lexical override dictionary $\mathcal{D}$, applied before general rules:

\begin{center}
\begin{tabular}{ll}
Latin & Arabic \\
\hline
i & \AR{آي} \\
na & \AR{نا} \\
no & \AR{نو} \\
dey & \AR{دي} \\
don & \AR{دون} \\
dem & \AR{ديم} \\
wetin & \AR{ويتين} \\
abeg & \AR{ابج} \\
\end{tabular}
\end{center}

This ensures that analytic particles remain visually discrete, which is central to the pedagogical function of the script.

\subsubsection{Ambiguity Management}

Because Arabic-only constraints collapse certain contrasts, two systematic ambiguities arise.

The first is the collapse of /g/ and /j/ under \AR{ج}. Resolution is lexical: in practice, Nigerian Pidgin has relatively few minimal pairs dependent on this contrast, and the intended reading is recoverable from context. If disambiguation is required, one may optionally reserve \AR{غ} for /g/ and \AR{ج} for /j/ as an internal convention, while still remaining within Arabic letters.

The second is the collapse of /p/ into \AR{ب} if the Latin baseline includes \texttt{p}. Under Arabic-only constraints, \texttt{p} is rewritten as \AR{ب}. This is historically normal in many Ajami systems. If one wishes to preserve /p/ without non-Arabic letters, the only option is a diacritic convention on \AR{ب}, but this is best treated as a later refinement rather than a core requirement.

\subsubsection{Algorithmic Order}

The transliteration pipeline is therefore:

First, apply lexical overrides $\mathcal{D}$ tokenwise. Second, apply normalization to remove English residue patterns. Third, apply greedy consonant mapping with longest-match priority. Fourth, apply the vowel strategy, either readability-first or pedagogical-diacritic mode. Finally, reattach punctuation and preserve spacing.

This order ensures that high-frequency particles stabilize early, consonantal structure is preserved deterministically, and vowels are added in a predictable manner.

\subsubsection{Worked Derivations}

\paragraph{Example 1:} \texttt{i dey}

Lexical override yields:
\[
\texttt{i} \mapsto \text{\AR{آي}},\qquad \texttt{dey} \mapsto \text{\AR{دي}}
\]
so the result is:
\[
\text{\AR{آي} \AR{دي}}.
\]

\paragraph{Example 2:} \texttt{wetin dey happen}

Lexical override yields:
\[
\texttt{wetin} \mapsto \text{\AR{ويتين}},\qquad \texttt{dey} \mapsto \text{\AR{دي}}.
\]
For \texttt{happen}, consonant mapping yields:
\[
h \mapsto \AR{ه},\quad p \mapsto \AR{ب},\quad n \mapsto \AR{ن}.
\]
Under readability-first vowels, a stable output is:
\[
\text{\AR{هابن}}.
\]
Therefore the full output is:
\[
\text{\AR{ويتين} \AR{دي} \AR{هابن؟}}
\]

\paragraph{Example 3:} \texttt{dem don come}

Lexical override yields:
\[
\texttt{dem} \mapsto \text{\AR{ديم}},\qquad \texttt{don} \mapsto \text{\AR{دون}}.
\]
For \texttt{come}, consonant mapping yields:
\[
c \mapsto \AR{ك},\quad m \mapsto \AR{م}.
\]
Vowels yield:
\[
o \mapsto \AR{و},\quad e \mapsto \AR{ي}.
\]
Thus:
\[
\texttt{come} \mapsto \text{\AR{كومي}}.
\]
A readability adjustment may drop the final \AR{ي} to reflect common pronunciation, producing:
\[
\text{\AR{كوم}}.
\]
The full output is:
\[
\text{\AR{ديم} \AR{دون} \AR{كوم}}.
\]

The derivations illustrate that the system is deterministic at the rule level while still permitting a controlled layer of phonological smoothing for high-frequency words.

\section{Comparison with Hausa Ajami}

Any proposal for an Arabic-based orthography for Nigerian Pidgin necessarily invites comparison with Hausa Ajami, the long-standing Arabic-script tradition used to write Hausa across northern Nigeria and the wider Sahel. The comparison is instructive because Hausa Ajami represents one of the most historically successful adaptations of Arabic script to a non-Arabic phonological system in West Africa.

\subsection{Historical Orientation}

Hausa Ajami emerged within Islamic scholarly networks and was used for poetry, legal documents, correspondence, and religious instruction. It developed organically over centuries, rather than through centralized planning. Orthographic conventions stabilized gradually through repeated manuscript transmission.

In contrast, the proposed Pidgin Ajami system is explicitly planned. It begins with a phonological inventory and a deterministic transliteration algorithm. The Hausa system evolved bottom-up; the Pidgin system is top-down.

\subsection{Phonological Adaptation}

Hausa contains phonemic contrasts not present in Arabic, including ejectives and implosives, as well as distinct /p/ and /g/ sounds. Hausa Ajami resolved these gaps in several ways:

First, some traditions introduced modified letters through dotting conventions.

Second, some contrasts were collapsed orthographically and disambiguated contextually.

Third, vowel length distinctions were often marked more carefully than in unvowelled Arabic.

The proposed Pidgin Ajami system aligns most closely with the second strategy. Under Arabic-only constraints, /p/ merges with \AR{ب} and /g/ may be represented by \AR{ج} or \AR{غ}. As in Hausa Ajami, phonological recovery is expected to occur through lexical familiarity rather than perfect graphemic fidelity.

The structural similarity lies in functional tolerance for controlled ambiguity.

\subsection{Vowel Representation}

Hausa distinguishes long and short vowels phonemically. In many Ajami manuscripts, long vowels are written explicitly with \AR{ا}, \AR{ي}, and \AR{و}, while short vowels may be omitted or marked with diacritics depending on genre.

Nigerian Pidgin does not rely heavily on vowel length contrasts for lexical differentiation. The proposed Pidgin Ajami therefore adopts a readability-first vowel strategy, writing core vowels explicitly while permitting optional diacritics for pedagogical precision.

Both systems therefore prioritize functional literacy over phonetic exhaustiveness.

\subsection{Morphological Transparency}

Hausa is morphologically richer than Nigerian Pidgin. It contains noun class remnants, plural patterns, and verbal derivations. Hausa Ajami thus encodes morphological alternations within stems.

Nigerian Pidgin, by contrast, is analytically structured. Tense, aspect, and modality are particle-based rather than inflectional. This produces a different orthographic opportunity.

In Pidgin Ajami, particles such as \AR{نا} (na), \AR{دي} (dey), \AR{دون} (don), and \AR{نو} (no) become visually stable units. Because these particles are invariant, the script reinforces analytic grammar rather than morphological alternation.

Hausa Ajami must accommodate internal stem alternation. Pidgin Ajami highlights syntactic particles.

\subsection{Sociolinguistic Position}

Hausa Ajami historically functioned alongside Arabic literacy and later alongside Latin-script Hausa. It did not replace Latin orthography but coexisted with it.

A Pidgin Ajami system would likely occupy a similar sociolinguistic niche: supplemental rather than replacement orthography. It could function in artistic, pedagogical, or cultural contexts without displacing established Latin usage.

The comparison therefore suggests that script plurality is not structurally destabilizing. West African linguistic ecology already supports dual-script traditions.

\subsection{Structural Divergence}

Despite parallels, there is a critical difference.

Hausa is an indigenous Niger–Congo language with a long-standing literary identity. Ajami writing reinforced an existing linguistic core.

Nigerian Pidgin is a contact language with English lexical ancestry and multi-ethnic substrate influence. Writing it in Arabic script visually decouples it from English orthography. This is not merely phonetic adaptation; it is perceptual reclassification.

Hausa Ajami preserves Hausa as Hausa. Pidgin Ajami repositions Pidgin as structurally autonomous.

\subsection{Implications for Standardization}

The Hausa case demonstrates that:

Orthographic conventions stabilize through repeated communal usage rather than abstract design alone.

Phoneme–grapheme mismatch does not prevent literacy if functional load remains manageable.

Script change does not erase Latin literacy but can coexist with it.

For Nigerian Pidgin, the lesson is not that a perfect mapping is required. Rather, consistency and repeatability are more important than exhaustiveness. A deterministic transliteration framework, once adopted in sustained textual practice, would likely converge toward stable norms just as Hausa Ajami did historically.

\section{Deterministic Reverse Transliteration into English}

This section formalizes the reverse mapping from the Arabic-based orthography into a normalized English phonetic baseline. The goal is structural symmetry: transliteration must be reversible at the phonemic level even if etymological English spelling is not perfectly reconstructed.

\subsection{Scope}

The reverse system maps Arabic grapheme sequences into a normalized English phonetic representation approximating contemporary pronunciation rather than historical spelling.

For example:

\[
\AR{كَابِتَالِسْتْ} \rightarrow kapitalist
\]

not necessarily the etymologically correct ``capitalist.''

\subsection{Preprocessing}

Let the input be a Unicode Arabic string. Processing proceeds tokenwise, separated by whitespace and punctuation.

Arabic diacritics (fatha, kasra, damma, sukun) are first stripped unless pedagogical vowel precision is required. This ensures consistent mapping independent of diacritic usage.

\subsection{Consonant Mapping}

The base consonant reversal map is:

\begin{center}
\begin{tabular}{lll}
Arabic & Latin & Notes \\
\hline
\AR{ب} & b \\
\AR{ت} & t \\
\AR{د} & d \\
\AR{ك} & k \\
\AR{ج} & g or j (contextual) \\
\AR{ف} & f \\
\AR{س} & s \\
\AR{ز} & z \\
\AR{ه} & h \\
\AR{م} & m \\
\AR{ن} & n \\
\AR{ر} & r \\
\AR{ل} & l \\
\AR{و} & w (consonant) \\
\AR{ي} & y (consonant) \\
\AR{ش} & sh \\
\AR{تش} & ch \\
\AR{غ} & g (if used as g marker) \\
\end{tabular}
\end{center}

Context rule for \AR{ج}:

If followed by front vowel marker (\AR{ي} or kasra), render as j.

Otherwise default to g.

This resolves most ambiguity without dictionary lookup.

\subsection{Vowel Mapping}

Long vowels:

\[
\AR{ا} \rightarrow a
\]
\[
\AR{ي} \rightarrow i
\]
\[
\AR{و} \rightarrow u
\]

When vowel letters appear between consonants, they are interpreted as full vowels rather than glides unless preceded by another vowel.

Example:

\[
\AR{رُولِنْزْ} \rightarrow rulins
\]

Optional smoothing rule:

u → oo when English-like representation is desired (e.g., rulins → rulins / rulings)

\subsection{Cluster Reconstruction}

Because Arabic orthography may collapse consonant clusters, a smoothing stage restores likely English clusters.

Example:

\[
\AR{سِمْبْتُمْزْ}
\]

Direct transliteration:

\[
simbtums
\]

Smoothing rule:

mbt → mpt

Thus:

\[
simptoms
\]

Final English baseline:

\[
symptoms
\]

This stage is heuristic and optional.

\subsection{Lexical Overrides}

As with forward transliteration, high-frequency forms may be stabilized via lexical overrides.

For example:

\[
\AR{ذَا} \rightarrow the
\]

\[
\AR{إِنْ} \rightarrow in
\]

\[
\AR{أُفْ} \rightarrow of
\]

\[
\AR{أَنْدْ} \rightarrow and
\]

\[
\AR{إِتْ} \rightarrow it
\]

These overrides allow restoration of conventional English articles and conjunctions.

\subsection{Worked Example}

Input:

\[
\AR{كَابِتَالِسْتْ} \AR{يُتُوبِيَانِيزْمْ}
\]

Step 1: Strip diacritics (optional).

\[
\AR{كابتالست} \AR{يتوبيانيزم}
\]

Step 2: Consonant map:

k a b t a l s t y t u b y a n i z m

Step 3: Vowel restoration:

kapitalist utopianizm

Step 4: English smoothing:

capitalist utopianism

Final output:

Capitalist Utopianism

\section*{Conclusion}

This study has proposed a structurally grounded introduction to Nigerian Pidgin through three interlocking lenses: grammatical description, phonological analysis, and orthographic engineering. Rather than treating the language as an informal derivative of English, the discussion has approached it as a system with its own internal regularities, aspectual structure, pronominal economy, and discourse pragmatics. The grammatical patterns examined demonstrate that Nigerian Pidgin is not a simplified variant of English but a rule-governed language whose apparent transparency conceals a coherent underlying architecture.

The orthographic proposal extends this structural orientation into the domain of script. By restricting the writing system to the Classical Arabic grapheme inventory and introducing minimal, explicitly defined digraph conventions, the system preserves typographic economy while maintaining phonological intelligibility. The result is neither an imitation of Arabic morphology nor an attempt to reproduce English spelling conventions. It is instead a phonologically motivated encoding scheme that reflects acoustic structure rather than etymology.

Crucially, the transliteration architecture has been formulated in explicitly compositional terms. The movement from Arabic graphemes to phonetic English and from phonetic English to standardized orthography demonstrates that transliteration is not a single act but a layered process. Each layer performs a distinct function: acoustic reconstruction, lexical normalization, and literary smoothing. By separating these stages formally, the system avoids conflating sound with spelling and representation with convention.

The broader significance of this framework lies in its methodological posture. A writing system need not be inherited to be legitimate; it may be engineered, provided its design principles are explicit. When phoneme inventories are clearly specified, grapheme correspondences are deterministically defined, and reading rules are computationally describable, orthography becomes a domain of formal design rather than historical accident.

Nigerian Pidgin, long transmitted primarily through speech and Latin-based informal writing, proves adaptable to such formalization. The exercise confirms that the language possesses sufficient structural stability to sustain cross-script representation without loss of grammatical identity. In doing so, it demonstrates that linguistic description, orthographic construction, and computational modeling need not be separate enterprises. They may converge in a single system grounded in phonology, constrained by script economy, and articulated through formal mapping.

What emerges is not merely a pedagogical experiment but a demonstration of principle: that structure precedes symbol, and that symbol, when responsibly engineered, can faithfully reflect the structure it encodes.

\newpage
\section*{Appendices}

\appendix
\section{Appendix A: Phoneme Inventory of Nigerian Pidgin}

\subsection{Methodological Note}

The phoneme inventory presented here reflects widely attested urban Nigerian Pidgin varieties. Variation exists regionally, particularly under substrate influence from Yoruba, Igbo, and Hausa; the inventory below represents a conservative shared core sufficient for orthographic modeling.

\subsection{Consonant Inventory}

The consonant phonemes may be organized by place and manner of articulation as follows.

\begin{center}
\begin{tabular}{lcccccc}
 & Bilabial & Labiodental & Alveolar & Postalveolar & Velar & Glottal \\ \hline
Plosive (voiceless) & p &  & t &  & k &  \\
Plosive (voiced) & b &  & d &  & g &  \\
Fricative (voiceless) &  & f & s & ʃ &  & h \\
Fricative (voiced) &  & v & z &  &  &  \\
Affricate &  &  &  & tʃ &  &  \\
Nasal & m &  & n &  & ŋ &  \\
Liquid &  &  & l, r &  &  &  \\
Glide & w &  &  &  & j &  \\
\end{tabular}
\end{center}

\subsection{Observations on Consonant Structure}

Consonant clusters are limited. Word-initial clusters are generally restricted to those inherited from English borrowings (e.g., \textit{school}, \textit{small}), though epenthesis or reduction frequently occurs in casual speech. Word-final obstruents are typically unreleased and may undergo devoicing.

The rhotic phoneme exhibits variable realization, ranging from alveolar tap to approximant depending on speaker background.

The velar nasal /ŋ/ occurs syllable-finally and in medial position but rarely word-initially.

\subsection{Vowel Inventory}

The vowel system is comparatively reduced relative to Standard English. A five- to seven-vowel analysis is common. A conservative six-vowel system is given below.

\begin{center}
\begin{tabular}{lccc}
 & Front & Central & Back \\ \hline
High & i &  & u \\
Mid & e &  & o \\
Low &  & a &  \\
\end{tabular}
\end{center}

Diphthongs such as /ai/, /au/, and /oi/ occur primarily in English-derived lexemes.

\subsection{Prosodic Structure}

Nigerian Pidgin is not tonally contrastive in the phonemic sense typical of many substrate languages. However, intonation patterns reflect substrate influence and may encode pragmatic distinctions.

Stress placement is typically inherited from English source forms but tends toward simplification in polysyllabic borrowings.

\subsection{Phonological Implications for Orthography}

The inventory reveals several structural constraints relevant to script design:

\begin{enumerate}
\item The vowel system is small and stable, making vowel marking feasible without excessive diacritic burden.
\item Consonant distinctions are largely preserved from English but without complex cluster productivity.
\item Absence of phonemic tone reduces orthographic complexity relative to tonal languages such as Yoruba.
\end{enumerate}

Any orthographic system must therefore preserve:

\[
\text{Contrast}(p,b),\;
\text{Contrast}(t,d),\;
\text{Contrast}(k,g),\;
\text{Contrast}(s,ʃ),\;
\text{Contrast}(f,v)
\]

while maintaining vowel distinctiveness sufficient to avoid lexical ambiguity.

\subsection{Conclusion}

This phoneme inventory provides the foundational layer upon which the proposed Arabic-script orthography is constructed. Orthographic evaluation must proceed downstream from phonological specification. Without this inventory, grapheme selection cannot be systematically justified.

\section{Appendix B: Grapheme--Phoneme Correspondence Table}

\subsection{Principles of Mapping}

The proposed orthography restricts itself to the Classical Arabic consonant inventory and uses a small number of explicit digraph conventions to represent Nigerian Pidgin phonemes not encoded as single Arabic letters. The aim is not perfect reversibility in the absence of vowel marks, but a stable pedagogical correspondence that can be made fully deterministic under full vowel marking.

Formally, the system defines a mapping
\[
\text{(phoneme or phoneme cluster)} \longrightarrow \text{(Arabic grapheme sequence)}
\]
with the convention that certain sequences of letters are treated as single units for reading purposes.

\subsection{Consonant Correspondence}

\begin{center}
\begin{tabular}{lll}
\textbf{Phoneme (IPA)} & \textbf{Arabic Grapheme(s)} & \textbf{Example (Pidgin)} \\ \hline
/p/ & \AR{ب} & \AR{بِيجِنْ} (pidgin) \\
/b/ & \AR{ب} & \AR{بُوِي} (boy) \\
/t/ & \AR{ت} & \AR{تِيكْ} (take) \\
/d/ & \AR{د} & \AR{دِمْ} (dem) \\
/k/ & \AR{ك} & \AR{كَمْ} (come) \\
/g/ & \AR{ك}\,(context) & \AR{كُو} (go) \\
/f/ & \AR{ف} & \AR{فِيتْ} (fit) \\
/v/ & \AR{ف} & \AR{فِرِي} (very) \\
/s/ & \AR{س} & \AR{سِي} (see) \\
/\textipa{\textesh}/ & \AR{ش} & \AR{شُورْ} (sure) \\
/z/ & \AR{ز} & \AR{زُو} (zoo) \\
/t\textipa{\textesh}/ & \AR{ت}\,+\,\AR{ش} & \AR{تْشُوبْ} (chop) \\
/m/ & \AR{م} & \AR{مِي} (me) \\
/n/ & \AR{ن} & \AR{نُو} (no) \\
/\textipa{N}/ & \AR{ن}\,(context) & \AR{سِنْ} (sing) \\
/l/ & \AR{ل} & \AR{لايْفْ} (life) \\
/r/ & \AR{ر} & \AR{رَنْ} (run) \\
/w/ & \AR{و} & \AR{وِي} (we) \\
/j/ & \AR{ي} & \AR{يَانْ} (yarn) \\
/h/ & \AR{ه} & \AR{هِيَرْ} (hear) \\
\end{tabular}
\end{center}

The system intentionally merges certain contrasts that are not represented in Classical Arabic as single letters. In particular, /p/ is written with \AR{ب} and /v/ is written with \AR{ف}. The affricate /t\textipa{\textesh}/ is written as the explicit digraph \AR{ت}+\AR{ش}, which is interpreted as a single consonant in reading.

\subsection{The \texorpdfstring{$\mathbf{\AR{ت}+\AR{ش}}$}{t+sh} Digraph Convention}

Let the digraph operator be defined by
\[
\mathrm{CH} := \text{\AR{ت}}\,\text{\AR{ش}}.
\]
Then the reading rule is:
\[
\mathrm{CH} \mapsto /t\textipa{\textesh}/
\]
regardless of whether the two letters could otherwise be parsed separately. When desired for pedagogical clarity, a suk\=un may be placed on \AR{ت} to signal clustering:
\[
\text{\AR{تْش}} \quad \text{rather than} \quad \text{\AR{تش}}.
\]
No additional letters beyond the Classical inventory are required.

\subsection{Vowel Correspondence}

The orthography permits full vowel marking in formal or instructional contexts and partial vowel omission in fluent contexts. A pedagogical one-to-one mapping is assumed when vowels are marked.

\begin{center}
\begin{tabular}{lll}
\textbf{Phoneme (IPA)} & \textbf{Arabic Representation} & \textbf{Example} \\ \hline
/i/ & \AR{ِ} (kasra) or \AR{ي} & \AR{سِي} (see) \\
/e/ & \AR{ِ} or \AR{َيْ} & \AR{بِلِيفْ} (believe) \\
/a/ & \AR{َ} (fatha) or \AR{ا} & \AR{مَاكْ} (make) \\
/o/ & \AR{ُ} (damma) or \AR{و} & \AR{كُو} (go) \\
/u/ & \AR{ُ} or \AR{و} & \AR{فُودْ} (food) \\
/ai/ & \AR{اي} & \AR{لايْفْ} (life) \\
/au/ & \AR{او} & \AR{ناو} (now) \\
/oi/ & \AR{اوي} & \AR{بُوِي} (boy) \\
\end{tabular}
\end{center}

\subsection{Orthographic Determinism Under Full Marking}

Let $\Sigma_P$ denote the intended Pidgin phoneme inventory and let $\Sigma_A^\ast$ denote finite Arabic grapheme strings. Define the encoding map
\[
G : \Sigma_P \rightarrow \Sigma_A^\ast
\]
with $G(t\textipa{\textesh})=\text{\AR{ت}}\text{\AR{ش}}$ as a primitive assignment. Under full vowel marking, the reading map becomes operationally deterministic for the intended learner audience. When vowels are omitted, ambiguity is treated as a controlled property of the script rather than a defect, paralleling ordinary Arabic orthographic practice.

\subsection{Cluster Representation}

Consonant clusters are written sequentially. Optional suk\=un marking may be used to prevent unintended vowel epenthesis in novice reading. No epenthetic vowels are introduced orthographically unless they are phonologically realized in the target utterance.

\subsection{Conclusion}

This correspondence table defines a Classical-Arabic-only Ajami-style orthography for Nigerian Pidgin with a minimal digraph extension for /t\textipa{\textesh}/. The result preserves typographic accessibility while keeping the mapping rules explicit and formally describable.

\section{Appendix C: Syllable Structure and Phonotactic Constraints}

\subsection{Syllable Template}

Nigerian Pidgin syllable structure may be described using the following general template:

\[
\sigma = (C)(C)V(C)
\]

where:

\begin{itemize}
\item \( C \) denotes a consonant,
\item \( V \) denotes a vowel nucleus,
\item parentheses indicate optional positions.
\end{itemize}

The most frequent syllable types are:

\[
CV,\quad CVC,\quad V,\quad CCV
\]

Complex codas are rare and typically inherited from English borrowings.

\subsection{Onset Structure}

Single-consonant onsets are unmarked:

\[
\text{go} \rightarrow CV
\]
\[
\text{fit} \rightarrow CVC
\]

Clustered onsets occur primarily in borrowings:

\[
\text{small} \rightarrow CCVC
\]
\[
\text{school} \rightarrow CCVC
\]

However, phonological reduction frequently occurs in rapid speech, yielding:

\[
\text{small} \rightarrow \text{[smɔl] or [s'mɔl]}
\]

There is no productive native cluster formation independent of English lexical inheritance.

\subsection{Coda Structure}

Word-final codas are typically single consonants:

\[
\text{fit} \rightarrow CVC
\]
\[
\text{run} \rightarrow CVC
\]

Complex codas such as /-nd/ or /-st/ are possible but are variably simplified:

\[
\text{last} \rightarrow [las] \text{ in casual speech}
\]

Thus coda clusters are lexically inherited rather than generatively productive.

\subsection{Vowel Reduction and Stability}

Unlike Standard English, vowel reduction in unstressed syllables is minimal. The system tends toward vowel preservation rather than centralized schwa collapse.

This relative vowel stability supports explicit vowel marking in orthography without excessive burden.

\subsection{Phonotactic Constraints}

The following constraints are broadly observed:

\[
\text{No word-initial } /ŋ/
\]

\[
\text{No triple consonant clusters}
\]

\[
\text{Limited coda complexity}
\]

\[
\text{Preference for open syllables in non-borrowed lexicon}
\]

\subsection{Implications for Arabic-Script Adaptation}

The phonotactic profile supports direct consonant sequencing without epenthetic insertion. Since clusters are limited and predictable, the orthography does not require additional structural markers beyond linear concatenation.

Let a Pidgin word be represented as:

\[
W = C_1 C_2 V C_3
\]

Orthographic representation preserves cluster integrity:

\[
C_1 C_2 V C_3 \rightarrow G(C_1)G(C_2)G(V)G(C_3)
\]

No morphophonemic alternations require special graphemic treatment.

\subsection{Formal Constraint Statement}

Define allowable syllable set \( \mathcal{S} \) such that:

\[
\mathcal{S} = \{ V, CV, CVC, CCV, CCVC \}
\]

with CCV and CCVC restricted to lexically inherited forms.

Orthographic validity requires:

\[
\forall \sigma \in \mathcal{S}, \; G(\sigma) \text{ is representable without additional diacritics beyond vowel marking.}
\]

\subsection{Conclusion}

The phonotactic system of Nigerian Pidgin is structurally simple relative to English and substrate languages. This simplicity facilitates orthographic determinism and minimizes ambiguity in Arabic-script adaptation.

\section{Appendix D: Morphosyntactic Structure and Functional Tagging}

\subsection{Overview}

This appendix provides a formal description of core morphosyntactic operators attested in the corpus. Nigerian Pidgin is largely analytic, with grammatical relations expressed through particles rather than inflectional morphology.

We adopt a functional notation in which lexical roots are represented as \( R \), grammatical particles as \( G \), and argument positions as \( A_i \).

\subsection{Tense–Aspect–Mood System}

The tense–aspect system is particle-based. Let \( V \) denote a verb root.

\subsubsection*{Progressive Marker \textit{dey}}

The progressive construction may be formalized as:

\[
S \rightarrow NP + \text{dey} + V
\]

Example:

\[
\text{I dey waka}
\]

\[
NP_1 + G_{PROG} + V
\]

Here, \textit{dey} encodes imperfective progressive aspect.

\subsubsection*{Perfective Marker \textit{don}}

The perfective construction may be formalized:

\[
S \rightarrow NP + \text{don} + V
\]

Example:

\[
\text{Dem don come}
\]

\[
NP_1 + G_{PERF} + V
\]

\textit{don} marks completed action without inflectional agreement.

\subsubsection*{Modal Marker \textit{fit}}

Ability or possibility is expressed via:

\[
S \rightarrow NP + \text{fit} + V
\]

Example:

\[
\text{You fit do am}
\]

This corresponds to modal \( CAN \).

\subsection{Negation}

Negation is pre-verbal and invariant:

\[
S \rightarrow NP + \text{no} + V
\]

Example:

\[
\text{I no sabi}
\]

There is no auxiliary inversion or do-support.

\subsection{Copular Constructions}

Two copular strategies exist.

\subsubsection*{Equative Copula \textit{na}}

\[
S \rightarrow \text{na} + NP
\]

Example:

\[
\text{Na true}
\]

\textit{na} marks equative identification.

\subsubsection*{Zero Copula}

Predicate adjectives often appear without overt copula:

\[
\text{She no well}
\]

This corresponds structurally to:

\[
NP + NEG + ADJ
\]

\subsection{Pronominal System}

The core pronoun inventory may be represented:

\begin{center}
\begin{tabular}{ll}
\textbf{Function} & \textbf{Form} \\ \hline
1SG & I \\
2SG & you \\
3SG & e / she / he \\
1PL & we \\
2PL & una \\
3PL & dem \\
\end{tabular}
\end{center}

Pronouns do not inflect for case.

\subsection{Object Marker \textit{am}}

Third-person objects frequently appear as invariant particle:

\[
S \rightarrow NP_1 + V + am
\]

Example:

\[
\text{I do am}
\]

This functions as a cliticized object pronoun.

\subsection{Serial Verb Construction}

Serial verb constructions occur without conjunction:

\[
NP + V_1 + V_2
\]

Example:

\[
\text{Make we go see}
\]

This reflects substrate influence from West African languages.

\subsection{Clause Typology}

Interrogatives are formed via intonation or question words without inversion:

\[
\text{You dey come?}
\]

No auxiliary movement occurs.

\subsection{Formal Summary}

Let:

\[
\mathcal{G} = \{\text{dey}, \text{don}, \text{fit}, \text{no}, \text{na}\}
\]

The grammar is characterized by:

\[
\text{No inflectional morphology}
\]

\[
\text{Pre-verbal TAM particles}
\]

\[
\text{Invariant negation}
\]

\[
\text{Minimal agreement marking}
\]

Thus Nigerian Pidgin may be modeled as an analytic grammar:

\[
S = NP + G^* + V + (NP)
\]

where \( G^* \) denotes zero or more grammatical particles.

\subsection{Implications for Orthography}

Because morphosyntactic distinctions are particle-based and not inflectional, orthographic clarity depends on stable representation of these particles:

\[
\text{dey}, \text{don}, \text{fit}, \text{no}, \text{na}
\]

Any ambiguity in their graphemic rendering would introduce grammatical ambiguity.

\subsection{Conclusion}

The morphosyntactic structure of Nigerian Pidgin is typologically analytic, particle-driven, and structurally transparent. Orthographic stability therefore depends less on morphological marking and more on consistent representation of invariant grammatical operators.

\section{Appendix E: Comparative Analysis with Hausa Ajami and West African Arabic-Script Traditions}

\subsection{Scope and Purpose}

This appendix situates the proposed Nigerian Pidgin Arabic-script system within the broader West African Ajami tradition. The objective is structural comparison rather than historical narrative.

Ajami refers to the adaptation of the Arabic script to represent non-Arabic languages. Hausa Ajami provides the most developed and standardized example in West Africa and therefore serves as the principal comparative baseline.

\subsection{Hausa Ajami: Structural Features}

Hausa Ajami adapts Arabic graphemes through three primary strategies:

\begin{enumerate}
\item Extension of grapheme inventory using modified letters (e.g., \AR{پ} for /p/, \AR{گ} for /g/).
\item Systematic marking of vowels via diacritics or long vowel letters.
\item Accommodation of tone indirectly through context rather than explicit marking.
\end{enumerate}

Hausa phonology includes implosives, ejectives, and tonal contrasts not present in Nigerian Pidgin. Consequently, Hausa Ajami employs additional graphemic innovation to encode these distinctions.

Formally, Hausa mapping may be expressed:

\[
G_H : \Sigma_{Hausa} \rightarrow \Sigma_{Ajami}
\]

with partial context sensitivity due to tone omission.

\subsection{Structural Comparison with Nigerian Pidgin}

Let:

\[
\Sigma_P = \text{Pidgin phoneme set}
\]
\[
\Sigma_H = \text{Hausa phoneme set}
\]

Then:

\[
|\Sigma_P| < |\Sigma_H|
\]

Pidgin lacks:

\[
\text{Tone contrast}
\]
\[
\text{Implosive consonants}
\]
\[
\text{Ejective consonants}
\]

Therefore, the orthographic burden required for Pidgin Ajami is strictly lower.

\subsection{Graphemic Extension Parallels}

Both Hausa Ajami and the proposed Pidgin system rely on extended Arabic letters:

\[
\AR{پ}, \AR{گ}, \AR{چ}, \AR{ڤ}
\]

These are not foreign innovations but established orthographic tools within Islamic manuscript traditions.

Thus the Pidgin proposal does not introduce structurally novel graphemes relative to regional precedent.

\subsection{Vowel Representation}

Hausa Ajami often under-specifies short vowels in non-pedagogical texts, relying on reader inference.

The proposed Pidgin system recommends optional full vowel marking in educational contexts and partial reduction in fluent writing.

This mirrors established Ajami practice:

\[
\text{Full marking} \rightarrow \text{pedagogical precision}
\]
\[
\text{Reduced marking} \rightarrow \text{scribal economy}
\]

\subsection{Morphological Transparency}

Hausa is morphologically richer than Pidgin. Hausa Ajami must encode:

\[
\text{Gender}
\]
\[
\text{Number agreement}
\]
\[
\text{Derivational morphology}
\]

Nigerian Pidgin, being analytic, imposes fewer morphological marking requirements.

Thus orthographic transparency may be achieved with lower complexity.

\subsection{Script Directionality and Cultural Integration}

Ajami scripts are written right-to-left. Adoption of Arabic script for Pidgin would therefore:

\[
\text{Align with established regional literacy traditions}
\]

However, unlike Hausa, Nigerian Pidgin has not historically been transmitted through Islamic scholarly networks. Therefore sociolinguistic adoption dynamics would differ.

\subsection{Formal Comparative Statement}

Let orthographic complexity be approximated as:

\[
C = |\Sigma_G| + D + T
\]

where:

\[
|\Sigma_G| = \text{number of graphemes}
\]
\[
D = \text{diacritic burden}
\]
\[
T = \text{tone marking complexity}
\]

For Hausa Ajami:

\[
T > 0
\]

For Nigerian Pidgin Ajami:

\[
T = 0
\]

Therefore:

\[
C_{Pidgin} < C_{Hausa}
\]

under equivalent vowel-marking regimes.

\subsection{Conclusion}

The proposed Pidgin Arabic orthography does not represent an unprecedented script innovation. It operates within established Ajami structural patterns while benefiting from reduced phonological complexity relative to Hausa.

The principal challenge is not graphemic feasibility but sociolinguistic adoption.

\section{Appendix F: Orthographic Economy and Information Density}

\subsection{Objective}

This appendix evaluates the proposed Arabic-script orthography for Nigerian Pidgin in terms of orthographic economy and information density. The aim is to quantify structural efficiency rather than to argue aesthetically.

We define orthographic economy as the ratio between phonological information encoded and graphemic length required.

\subsection{Definitions}

Let:

\[
W_L = \text{Latin-script representation length}
\]

\[
W_A = \text{Arabic-script representation length}
\]

Let \( |W| \) denote grapheme count excluding whitespace.

Define compression ratio:

\[
R = \frac{|W_A|}{|W_L|}
\]

\subsection{Empirical Sample}

Consider representative examples from the corpus.

\[
\text{I dey come}
\]

Latin length:

\[
|W_L| = 9
\]

Arabic:

\[
\AR{آيْ} \AR{دي} \AR{كَمْ}
\]

\[
|W_A| = 8
\]

Thus:

\[
R = \frac{8}{9} \approx 0.89
\]

Second example:

\[
\text{You no get sense}
\]

Latin:

\[
|W_L| = 14
\]

Arabic:

\[
\AR{يُو} \AR{نَوْ} \AR{جِتْ} \AR{سِنْسْ}
\]

\[
|W_A| = 13
\]

\[
R \approx 0.93
\]

Across sampled entries, preliminary averaging yields:

\[
R \in [0.85, 1.05]
\]

indicating near parity in graphemic density.

\subsection{Vowel Marking and Density}

If full diacritic marking is retained, character count increases marginally.

If short vowels are omitted in fluent writing, the ratio decreases:

\[
R_{reduced} < R_{full}
\]

Thus vowel omission provides optional compression at the cost of disambiguation reliance.

\subsection{Information-Theoretic Framing}

Let:

\[
H(P) = -\sum p(x)\log p(x)
\]

represent entropy of phoneme distribution.

Orthographic efficiency may be approximated as:

\[
E = \frac{H(P)}{\text{Average grapheme length}}
\]

Since Pidgin has relatively low phonemic entropy due to reduced vowel inventory and limited cluster productivity, the system supports compact encoding.

\subsection{Redundancy Analysis}

Latin orthography for Pidgin inherits English irregularities, including silent letters and digraph ambiguity:

\[
\text{“ough” phenomena absent but analogous ambiguity exists}
\]

Arabic-script representation is phonemic and thus reduces redundancy.

Redundancy \( \rho \) may be approximated as:

\[
\rho = 1 - \frac{\text{phoneme count}}{\text{grapheme count}}
\]

The proposed system aims to minimize \( \rho \).

\subsection{Comparative Economy}

Let:

\[
C_L = \text{Latin complexity}
\]
\[
C_A = \text{Arabic complexity}
\]

If Latin representation is quasi-etymological while Arabic representation is phonemic, then:

\[
C_A \leq C_L
\]

under consistent vowel marking conventions.

\subsection{Cognitive Load Considerations}

Right-to-left directionality introduces initial learning cost. However, grapheme inventory size remains modest:

\[
|\Sigma_A| \approx 29
\]

which lies within standard literacy acquisition norms.

\subsection{Conclusion}

Quantitative evaluation suggests that the proposed orthography:

\begin{enumerate}
\item Maintains near parity in graphemic length.
\item Reduces phoneme-to-grapheme ambiguity.
\item Allows optional compression through diacritic omission.
\end{enumerate}

From an information-density perspective, the system is structurally efficient and does not introduce measurable expansion relative to Latin representation.

\section{Appendix G: Ambiguity and Error Propagation Analysis}

\subsection{Objective}

This appendix evaluates potential sources of ambiguity and error propagation within the proposed Arabic-script orthography for Nigerian Pidgin. The goal is to identify structural weaknesses and determine whether they arise from the phonology of the language or from script adaptation.

\subsection{Formal Model of Orthographic Mapping}

Let:

\[
G : \Sigma_P^* \rightarrow \Sigma_A^*
\]

where \( \Sigma_P \) is the phoneme inventory and \( \Sigma_A \) the grapheme inventory.

Ambiguity arises if:

\[
\exists x,y \in \Sigma_P^*, \; x \neq y \;\text{such that}\; G(x) = G(y)
\]

We analyze potential cases.

\subsection{Vowel Omission Ambiguity}

If short vowels are omitted, sequences such as:

\[
\text{fit} \quad \text{and} \quad \text{fat}
\]

could reduce to identical consonantal skeletons.

Example:

\[
\AR{فْتْ}
\]

Thus vowel omission introduces recoverable but real ambiguity.

Mitigation strategy: mandatory vowel marking in minimal-pair environments.

\subsection{Homophony Inherited from Pidgin}

Certain lexical items are homophonous independent of script.

For example:

\[
\text{no} \quad (\text{negation}) 
\]
\[
\text{know} \quad (\text{lexical verb})
\]

In Pidgin these collapse phonologically:

\[
\AR{نو}
\]

The ambiguity is phonemic, not orthographic. Therefore the script does not introduce additional instability beyond the spoken language.

\subsection{Consonant Cluster Compression}

Clusters such as:

\[
\AR{سْمُولْ}
\]

may be visually dense but do not collapse phonemic distinction. There is no graphemic ambiguity provided cluster letters remain distinct.

\subsection{Loanword Adaptation}

English-derived terms may exhibit variable vowel realization. For example:

\[
\text{school} \rightarrow \AR{سكول}
\]

If speakers variably pronounce epenthetic vowels, orthographic normalization may diverge from speech variation. This represents standardization tension rather than mapping ambiguity.

\subsection{Diacritic Loss in Informal Writing}

If diacritics are omitted entirely:

\[
\AR{كتب}
\]

Ambiguity increases exponentially.

Define diacritic omission function:

\[
D : \Sigma_A \rightarrow \Sigma_A'
\]

where \( \Sigma_A' \subset \Sigma_A \) excludes short vowel marks.

Then ambiguity probability increases as:

\[
P(\text{collision}) \propto \frac{|\Sigma_P|}{|\Sigma_A'|}
\]

Full vowel marking therefore reduces collision probability.

\subsection{Error Propagation in Iterative Copying}

Consider iterative transcription:

\[
W_0 \rightarrow W_1 \rightarrow \dots \rightarrow W_n
\]

If each stage introduces independent error probability \( \epsilon \), then expected distortion after \( n \) iterations is:

\[
E[D_n] = 1 - (1-\epsilon)^n
\]

Phonemic orthography reduces \( \epsilon \) relative to quasi-etymological systems, since fewer silent letters or irregular mappings exist.

\subsection{Structural Stability Conditions}

The orthography remains injective under the following constraints:

\[
\text{Full vowel marking}
\]
\[
\text{Retention of extended graphemes}
\]
\[
\text{No grapheme conflation}
\]

If these conditions hold, then:

\[
G \text{ is injective}
\]

If diacritics are omitted:

\[
G \text{ becomes partially many-to-one}
\]

\subsection{Comparative Ambiguity}

Latin-script Pidgin already contains ambiguity due to English spelling conventions. For example:

\[
\text{read} \; (\text{present}) \neq \text{read} \; (\text{past})
\]

Such orthographic instability is inherited from English rather than intrinsic to Pidgin.

The Arabic-script proposal, being phonemic, may reduce certain inherited ambiguities.

\subsection{Conclusion}

Ambiguity in the proposed system arises primarily from optional vowel omission rather than from consonant representation. When vowels are fully marked, the mapping is structurally robust and minimally ambiguous.

Error propagation rates are comparable to, and potentially lower than, Latin-script representation due to phonemic transparency.

\section{Appendix H: Reversibility Theorem and Transliteration Functions}

\subsection{Definitions}

Let:

\[
\Sigma_P = \text{Pidgin phoneme inventory}
\]

\[
\Sigma_A = \text{Arabic grapheme inventory}
\]

\[
\Sigma_L = \text{Latin grapheme inventory}
\]

Define transliteration functions:

\[
G : \Sigma_P^* \rightarrow \Sigma_A^*
\]

\[
L : \Sigma_P^* \rightarrow \Sigma_L^*
\]

and inverse mappings:

\[
G^{-1} : \Sigma_A^* \rightarrow \Sigma_P^*
\]

\[
L^{-1} : \Sigma_L^* \rightarrow \Sigma_P^*
\]

\subsection{Orthographic Determinism}

The proposed system assigns a unique grapheme to each phoneme:

\[
\forall p \in \Sigma_P,\; G(p) = a \in \Sigma_A
\]

such that:

\[
p_1 \neq p_2 \Rightarrow G(p_1) \neq G(p_2)
\]

Thus:

\[
G \text{ is injective at the phoneme level.}
\]

\subsection{Weak Reversibility Theorem}

\textbf{Theorem.}  
If all vowel diacritics are preserved, then:

\[
G^{-1} \circ G = \mathrm{Id}_{\Sigma_P^*}
\]

That is, transliteration to Arabic script and back to phonemic representation recovers the original sequence.

\textbf{Proof.}

Since \( G \) is injective over \( \Sigma_P \), and since concatenation preserves symbol boundaries, we have:

\[
G(p_1 p_2 \dots p_n) = G(p_1)G(p_2)\dots G(p_n)
\]

Because each \( G(p_i) \) corresponds uniquely to \( p_i \), the inverse mapping recovers:

\[
G^{-1}(G(p_i)) = p_i
\]

Thus:

\[
G^{-1}(G(p_1 p_2 \dots p_n)) = p_1 p_2 \dots p_n
\]

provided no grapheme ambiguity is introduced by diacritic omission.

\hfill \( \square \)

\subsection{Partial Reversibility Under Vowel Omission}

Let \( D \) be a diacritic removal function:

\[
D : \Sigma_A^* \rightarrow \Sigma_A'^*
\]

where \( \Sigma_A' \subset \Sigma_A \) excludes short vowel markers.

Then composite mapping:

\[
G' = D \circ G
\]

may no longer be injective.

\textbf{Corollary.}  
If two phonemic sequences differ only in vowel quality and short vowels are unmarked, then:

\[
G'(x) = G'(y)
\]

even if \( x \neq y \).

Thus reversibility becomes context-dependent rather than absolute.

\subsection{Latin Transliteration Comparison}

Latin-script Pidgin is not strictly phonemic. Let:

\[
L(p) = \ell
\]

English orthographic conventions introduce non-phonemic irregularities. Therefore:

\[
L^{-1} \text{ is not strictly deterministic without phonological knowledge.}
\]

Hence, under full vowel marking:

\[
\text{Arabic-script representation may exhibit stronger phonemic reversibility than Latin representation.}
\]

\subsection{Algorithmic Implementation}

Define transliteration procedure:

\[
T_A: \text{Input Latin Pidgin} \rightarrow \text{Phonemic parse} \rightarrow G(\cdot)
\]

\[
T_L: \text{Input Arabic Pidgin} \rightarrow G^{-1}(\cdot) \rightarrow \text{Latin rendering}
\]

Since the mapping operates at phoneme level rather than orthographic level, intermediate phonemic parsing ensures structural consistency.

\subsection{Complexity}

Let word length be \( n \).

Then transliteration complexity:

\[
O(n)
\]

as each symbol is mapped independently.

No recursive or context-sensitive rewriting is required under full marking.

\subsection{Conclusion}

The proposed orthography satisfies:

\[
G^{-1} \circ G = \mathrm{Id}
\]

under full vowel marking.

Under vowel reduction, reversibility is weakened but remains recoverable via lexical context.

Thus the system is formally well-defined, computationally tractable, and structurally reversible within specified constraints.

\section{Appendix I: Minimal Pair Dataset and Contrast Preservation}

\subsection{Objective}

This appendix evaluates whether the proposed Arabic-script orthography preserves phonemic contrasts attested in Nigerian Pidgin. Orthographic adequacy requires that minimal phonemic distinctions correspond to distinct graphemic forms.

Formally, if:

\[
x, y \in \Sigma_P^*, \quad x \neq y
\]

and

\[
\text{Contrast}(x,y) \text{ is phonemic}
\]

then it must hold that:

\[
G(x) \neq G(y)
\]

\subsection{Consonant Contrast Tests}

\subsubsection*{Plosive Voicing Contrast}

\[
\text{pat} \neq \text{bat}
\]

\[
\AR{پَتْ} \neq \AR{بَتْ}
\]

Distinct graphemes preserve contrast.

\subsubsection*{Velar Contrast}

\[
\text{go} \neq \text{ko}
\]

\[
\AR{گُو} \neq \AR{كُو}
\]

Contrast preserved.

\subsubsection*{Fricative Contrast}

\[
\text{fan} \neq \text{van}
\]

\[
\AR{فَن} \neq \AR{ڤَن}
\]

Distinct graphemes maintain phonemic distinction.

\subsubsection*{Sibilant Contrast}

\[
\text{see} \neq \text{she}
\]

\[
\AR{سِي} \neq \AR{شِي}
\]

Contrast preserved.

\subsection{Vowel Contrast Tests}

\subsubsection*{/i/ vs /a/}

\[
\text{fit} \neq \text{fat}
\]

\[
\AR{فِيتْ} \neq \AR{فَتْ}
\]

Distinct vowel marking required for full preservation.

\subsubsection*{/o/ vs /u/}

\[
\text{cot} \neq \text{cut}
\]

\[
\AR{كُتْ} \neq \AR{كَتْ}
\]

Vowel marking disambiguates.

\subsection{Modal vs Lexical Distinction}

\[
\text{fit} \; (\text{modal}) \neq \text{fit} \; (\text{adjective})
\]

These are homophonous in Pidgin and therefore not orthographically distinguishable in any phonemic system.

Thus:

\[
\text{Ambiguity inherited from phonology}
\]

not from script.

\subsection{Particle Contrast Preservation}

\[
\text{no} \neq \text{don}
\]

\[
\AR{نو} \neq \AR{دُونْ}
\]

Distinct grammatical operators remain visually distinct.

\subsection{Cluster Preservation}

\[
\text{small} \neq \text{mall}
\]

\[
\AR{سْمُولْ} \neq \AR{مُولْ}
\]

Initial cluster preserved without conflation.

\subsection{Formal Statement}

Let:

\[
\mathcal{M} = \{(x,y) \mid x,y \text{ minimal pairs}\}
\]

The orthography satisfies:

\[
\forall (x,y) \in \mathcal{M}, \quad G(x) \neq G(y)
\]

under full vowel marking.

If vowel marking is omitted:

\[
\exists (x,y) \in \mathcal{M} \text{ such that } G'(x) = G'(y)
\]

only when the contrast depends solely on short vowels.

\subsection{Contrast Preservation Theorem}

\textbf{Theorem.}  
Under full vowel marking, the proposed orthography preserves all consonantal and vocalic phonemic contrasts present in the defined phoneme inventory.

\textbf{Proof Sketch.}

Since mapping \( G \) is injective at phoneme level and since minimal pairs differ in at least one phoneme, their grapheme sequences must differ in at least one position.

\[
x \neq y \Rightarrow \exists i \; (p_i \neq q_i)
\]

\[
G(p_i) \neq G(q_i)
\]

Thus:

\[
G(x) \neq G(y)
\]

\hfill \( \square \)

\subsection{Conclusion}

The orthography preserves phonemic contrast under full specification. Ambiguities arise only under deliberate vowel omission and are therefore optional rather than structural failures.

\section{Appendix J: Computational Transliteration Specification Under Classical Arabic Constraints}

\subsection{Alphabet Restriction}

The Arabic grapheme inventory is restricted to:

\[
\Sigma_A^{classical}
=
\{ \AR{ا}, \AR{ب}, \AR{ت}, \AR{ث}, \AR{ج}, \AR{ح}, \AR{خ}, \AR{د}, \AR{ذ}, \AR{ر}, \AR{ز}, \AR{س}, \AR{ش}, \AR{ص}, \AR{ض}, \AR{ط}, \AR{ظ}, \AR{ع}, \AR{غ}, \AR{ف}, \AR{ق}, \AR{ك}, \AR{ل}, \AR{م}, \AR{ن}, \AR{ه}, \AR{و}, \AR{ي} \}
\]

No Persian-derived letters (\AR{پ}, \AR{گ}, \AR{چ}, \AR{ڤ}) are permitted.

\subsection{Phoneme Set}

Let:

\[
\Sigma_P = \text{Pidgin phoneme inventory defined in Appendix A}
\]

including phonemes absent in Classical Arabic:

\[
\{ p, g, v, tʃ \}
\]

\subsection{Digraph Encoding Strategy}

Since injective single-grapheme mapping is impossible for these phonemes, we define a sequence-level mapping:

\[
G : \Sigma_P \rightarrow (\Sigma_A^{classical})^+
\]

where the codomain permits multi-grapheme outputs.

The following deterministic digraph encodings are adopted:

\[
p \rightarrow \AR{بْه}
\]

\[
g \rightarrow \AR{كْ}
\]

\[
v \rightarrow \AR{فْ}
\]

\[
tʃ \rightarrow \AR{تْش}
\]

These sequences are chosen to satisfy:

\[
G(p) \neq G(b)
\]

\[
G(g) \neq G(k)
\]

\[
G(v) \neq G(f)
\]

\[
G(tʃ) \neq G(t), G(ʃ)
\]

Thus injectivity is preserved at the sequence level.

\subsection{Latin-to-Phoneme Parsing}

Let:

\[
\pi : \Sigma_L^* \rightarrow \Sigma_P^*
\]

Parsing rules include maximal matching:

\[
\text{sh} \rightarrow ʃ
\]

\[
\text{ch} \rightarrow tʃ
\]

\[
\text{ng} \rightarrow ŋ
\]

\[
\text{ai} \rightarrow ai
\]

\[
\text{au} \rightarrow au
\]

Singleton mappings apply otherwise.

\subsection{Phoneme-to-Arabic Rewrite Function}

For each \( p \in \Sigma_P \):

\[
p \rightarrow G(p)
\]

Core mappings:

\[
b \rightarrow \AR{ب}
\]

\[
t \rightarrow \AR{ت}
\]

\[
d \rightarrow \AR{د}
\]

\[
k \rightarrow \AR{ك}
\]

\[
f \rightarrow \AR{ف}
\]

\[
ʃ \rightarrow \AR{ش}
\]

\[
m \rightarrow \AR{م}
\]

\[
n \rightarrow \AR{ن}
\]

\[
l \rightarrow \AR{ل}
\]

\[
r \rightarrow \AR{ر}
\]

\[
w \rightarrow \AR{و}
\]

\[
j \rightarrow \AR{ي}
\]

\[
h \rightarrow \AR{ه}
\]

Extended phonemes use digraph encoding defined above.

\subsection{Injectivity Theorem Under Digraph Encoding}

\textbf{Theorem.}  
Under the digraph encoding strategy, the transliteration mapping

\[
G : \Sigma_P^* \rightarrow (\Sigma_A^{classical})^*
\]

is injective.

\textbf{Proof.}

Assume:

\[
x \neq y
\]

Then there exists position \( i \) such that:

\[
p_i \neq q_i
\]

Since:

\[
G(p_i) \neq G(q_i)
\]

either as single grapheme or digraph sequence, concatenation preserves positional distinction.

Thus:

\[
G(x) \neq G(y)
\]

\hfill \( \square \)

\subsection{Reverse Mapping}

Define inverse mapping:

\[
G^{-1} : (\Sigma_A^{classical})^* \rightarrow \Sigma_P^*
\]

Parsing rule:

\begin{enumerate}
\item Detect digraph sequences first:
\[
\AR{بْه} \rightarrow p
\]
\[
\AR{كْ} \rightarrow g
\]
\[
\AR{فْ} \rightarrow v
\]
\[
\AR{تْش} \rightarrow tʃ
\]
\item Otherwise apply single-character mapping.
\end{enumerate}

Maximal-sequence matching ensures deterministic parsing.

\subsection{Algorithmic Implementation}

The system may be implemented as deterministic finite-state transducer:

\[
M = (Q, \Sigma_L, \Sigma_A^{classical}, \delta, q_0, F)
\]

with priority rules for digraph detection.

Time complexity remains:

\[
O(n)
\]

for input length \( n \).

\subsection{Correctness Condition}

Under full vowel marking and digraph preservation:

\[
T_L(T_A(w)) = w
\]

for all phonemically spelled Latin Pidgin input \( w \).

\subsection{Structural Trade-Off}

The removal of Persian letters increases average grapheme length:

\[
|G(p)| \ge 2
\]

for certain phonemes.

However, it preserves:

\[
\text{Injectivity}
\]
\[
\text{Finite-state computability}
\]
\[
\text{Reversibility}
\]

while maintaining strict adherence to Classical Arabic grapheme inventory.

\subsection{Conclusion}

Even under the constraint of Classical Arabic letters only, the transliteration system remains:

\[
\text{Deterministic}
\]
\[
\text{Injective}
\]
\[
\text{Computationally tractable}
\]

The cost is increased graphemic length for non-Arabic phonemes, but structural integrity is preserved.

\section{Appendix K: Arabic-Script Passage and Reverse Transliteration}

\subsection{Original Arabic-Script Text}

\begin{quote}
\AR{كَابِتَالِسْتْ} \AR{يُتُوبِيَانِيزْمْ}: \AR{ذَا} \AR{رِبْرِشَنْ} \AR{أُفْ} \AR{لَاكْ} \AR{أَنْدْ} \AR{إِتْسْ} \AR{كَلْتْشَرَلْ} \AR{سِمْبْتُمْزْ}

\AR{إِنْ} \AR{ذِسْ} \AR{بَارْتْ} \AR{أُفْ} \AR{هِيلِنْ} \AR{رُولِنْزْ} \AR{سَايْكُوسِينِمَا،} \AR{شِي} \AR{دِلْفْزْ} \AR{إِنْتُو} \AR{كَابِتَالِسْتْ} \AR{آيدِيُولُوجِي} \AR{أَزْ} \AR{إِتْ} \AR{إِنْتِرْسِكْتْسْ} \AR{وِذْ} \AR{ذَا} \AR{سَايْكُوأَنَالِيتِكْ} \AR{كُونْسِبْتْ} \AR{أُفْ} "\AR{لَاكْ}".

\AR{شِي} \AR{فْرَيْمْزْ} \AR{كَابِتَالِيزْمْ} \AR{أَزْ} \AR{أَ} \AR{سِسْتَمْ} \AR{بِلْتْ} \AR{أُونْ} \AR{دِنَايَلْ} \AR{أُفْ} \AR{ذِسْ} \AR{لَاكْ،} \AR{رِزَلْتِنْ} \AR{إِنْ} \AR{سَايْكُولُوجِيكَلْ} \AR{أَنْدْ} \AR{سُوشَلْ} \AR{كُنْسِكْوِنْسِزْ}.
\end{quote}

\subsection{Direct Phonetic Transliteration}

Capitalist Utopianism: The Repression of Lack and Its Cultural Symptoms

In this part of Heilin Rulins Psychocinema, she delves into capitalist ideology as it intersects with the psychoanalytic concept of “lack.”

She frames capitalism as a system built on denial of this lack, resulting in psychological and social consequences.

\subsection{Standardized English Rendering}

Capitalist Utopianism: The Repression of Lack and Its Cultural Symptoms

In this section of Helen Rollins’ *Psychocinema*, she explores capitalist ideology as it intersects with the psychoanalytic concept of “lack,” understood as an internal gap within the subject.

She argues that capitalism is structured around denial of this constitutive lack, producing both psychological and social consequences.

\subsection{Layered Observations}

The Arabic-script version encodes phonetic English rather than etymological English. Reverse transliteration first yields a phonetic baseline. A secondary smoothing stage restores conventional spelling and standard literary form.

The three layers may be schematized as:

\[
\text{Arabic grapheme} \rightarrow
\text{phonetic English} \rightarrow
\text{standardized English}.
\]

This layered structure confirms that the transliteration system is acoustically grounded while still allowing conventional normalization when required.
\section{Appendix L: Core Nigerian Pidgin Lexicon in Arabic and Latin Scripts}

\begin{longtable}{p{3.2cm} p{3.2cm} p{3.2cm} p{3.2cm}}
\textbf{Pidgin (Arabic)} & \textbf{Pidgin (Latin)} & \textbf{English (Arabic)} & \textbf{English (Latin)} \\ \hline
\endfirsthead
\textbf{Pidgin (Arabic)} & \textbf{Pidgin (Latin)} & \textbf{English (Arabic)} & \textbf{English (Latin)} \\ \hline
\endhead

\AR{هَوْ} \AR{فَرْ} \AR{ناو؟} & How far now? & \AR{هَوْ} \AR{آرْ} \AR{يُو؟} & How are you? \\
\AR{آيْ} \AR{دي} & I dey & \AR{آيْ} \AR{أَمْ} \AR{هِيَرْ} & I am here \\
\AR{آيْ} \AR{نو} \AR{نو} & I no know & \AR{آيْ} \AR{دُونْتْ} \AR{نُو} & I don't know \\
\AR{وِيتِنْ} \AR{دِي} \AR{آبُن؟} & Wetin dey happen? & \AR{وَتْ} \AR{إِزْ} \AR{جُوِينْ} \AR{آنْ؟} & What is going on? \\
\AR{نا} \AR{واو} & Na wa & \AR{ذَتْسْ} \AR{كْرَيْزِي} & That’s crazy \\
\AR{آيْ} \AR{بِلِيفْ} \AR{سِي} & I believe say & \AR{آيْ} \AR{ثِنْكْ} \AR{ذَتْ} & I think that \\
\AR{يُو} \AR{دي} \AR{كَمْ؟} & You dey come? & \AR{آرْ} \AR{يُو} \AR{كَمِينْجْ؟} & Are you coming? \\
\AR{نُو} \AR{واهلا} & No wahala & \AR{نُو} \AR{بْرُوبْلِمْ} & No problem \\
\AR{وِايْ} \AR{ناو؟} & Why now? & \AR{وِايْ} \AR{آرْ} \AR{يُو} \AR{دُوِينْ} \AR{ذِسْ؟} & Why are you doing this? \\
\AR{آيْ} \AR{دي} \AR{كَمْ} & I dey come & \AR{آيْ} \AR{آمْ} \AR{آنْ} \AR{مَايْ} \AR{وِايْ} & I’m on my way \\
\AR{نُو} \AR{مينْ} & No mean & \AR{آيْ} \AR{دِدْنْتْ} \AR{مِينْ} \AR{إِتْ} & I didn’t mean it \\
\AR{وِايْ} \AR{يُو} \AR{دُو} \AR{آمْ؟} & Why you do am? & \AR{وِايْ} \AR{دِدْ} \AR{يُو} \AR{دُو} \AR{إِتْ؟} & Why did you do it? \\
\AR{شِي} \AR{نو} \AR{وِيلْ} & She no well & \AR{شِي} \AR{إِزْ} \AR{نَاتْ} \AR{فِيلِينْجْ} \AR{وِيلْ} & She is not feeling well \\
\AR{دِمْ} \AR{دُونْ} \AR{كَمْ} & Dem don come & \AR{ذَيْ} \AR{آرْ} \AR{هِيَرْ} & They are here \\
\AR{بُوِي} \AR{دِي} \AR{كَمْ} & Boy dey come & \AR{ذَ} \AR{بُوِي} \AR{إِزْ} \AR{كَمِينْجْ} & The boy is coming \\
\AR{آيْ} \AR{دِي} \AR{هُنْجَرْ} & I dey hunger & \AR{آيْ} \AR{أَمْ} \AR{هَنْجْرِي} & I’m hungry \\
\AR{يُو} \AR{سَابِي} \AR{بوكْ؟} & You sabi book? & \AR{دُو} \AR{يُو} \AR{نُو} \AR{أَنِيثِينْجْ؟} & Do you know anything? \\
\AR{وِيتِنْ} \AR{بِي} \AR{دِسْ؟} & Wetin be this? & \AR{وَتْ} \AR{إِزْ} \AR{ذِسْ؟} & What is this? \\
\AR{آيْ} \AR{نو} \AR{فِيتْ} & I no fit & \AR{آيْ} \AR{كَانْ} \AR{نَاتْ} & I can’t \\
\AR{آيْ} \AR{دِي} \AR{تَايَرْ} & I dey tire & \AR{آيْ} \AR{أَمْ} \AR{تَايَرْدْ} & I am tired \\

\AR{أَبِجْ} & Abeg & \AR{بْلِيزْ} & Please \\
\AR{وَاها} \AR{لا} \AR{دِي} & Wahala dey & \AR{ذِيْرْ} \AR{إِزْ} \AR{بْرُوبْلِمْ} & There is a problem \\
\AR{نُو} \AR{وَاها} \AR{لا} \AR{دِي} & No wahala dey & \AR{نُو} \AR{بْرُوبْلِمْ} & No problem \\
\AR{يُو} \AR{دُونْ} \AR{شُورْ} & You don sure & \AR{آرْ} \AR{يُو} \AR{شُورْ؟} & Are you sure? \\
\AR{آيْ} \AR{نو} \AR{جِيتْ} & I no get & \AR{آيْ} \AR{دُونْتْ} \AR{هَافْ} & I don’t have \\
\AR{دِمْ} \AR{دَي} \AR{بْلَي} & Dem dey play & \AR{ذَيْ} \AR{أَرْ} \AR{نَاتْ} \AR{سِيرِيُوسْ} & They are not serious \\
\AR{مَاكْ} \AR{وِي} \AR{جُو} & Make we go & \AR{لِتْسْ} \AR{جُو} & Let’s go \\
\AR{كَالْمْ} \AR{داوْنْ} & Calm down & \AR{رِلَاكْسْ} & Relax \\
\AR{يُو} \AR{سَابِي} \AR{آمْ؟} & You sabi am? & \AR{دُو} \AR{يُو} \AR{نُو} \AR{إِتْ؟} & Do you know it? \\
\AR{آيْ} \AR{نَوْ} \AR{دِي} & I no dey & \AR{آيْ} \AR{أَمْ} \AR{أَوْكَيْ} & I’m okay \\
\AR{نَا} \AR{تْرُو} & Na true & \AR{ذَتْ} \AR{إِزْ} \AR{تُرُو} & That is true \\
\AR{نَا} \AR{لايْ} & Na lie & \AR{ذَتْ} \AR{إِزْ} \AR{أَ} \AR{لايْ} & That is a lie \\
\AR{يُو} \AR{دُونْ} \AR{تْرَاي} & You don try & \AR{يُو} \AR{دِدْ} \AR{وِلْ} & You did well \\
\AR{آيْ} \AR{نَوْ} \AR{فِيتْ} \AR{دُو} \AR{آمْ} & I no fit do am & \AR{آيْ} \AR{كَانْتْ} \AR{دُو} \AR{إِتْ} & I can’t do it \\
\AR{دِسْ} \AR{تِنْ} \AR{سْوِيتْ} & This thing sweet & \AR{ذِسْ} \AR{إِزْ} \AR{جُودْ} & This is good \\
\AR{آيْ} \AR{دِي} \AR{وَاكَا} & I dey waka & \AR{آيْ} \AR{أَمْ} \AR{وُوكِنْجْ} & I am walking \\
\AR{يُو} \AR{دِي} \AR{شَوْ} & You dey show & \AR{يُو} \AR{آرْ} \AR{بْرَاجِنْجْ} & You are bragging \\
\AR{نَا} \AR{سُوْ} \AR{سُوْ} & Na so so & \AR{أَلْوَيْزْ} / \AR{تُو} \AR{مُوشْ} & Always / too much \\
\AR{آيْ} \AR{دُونْ} \AR{تَايَرْ} & I don tire & \AR{آيْ} \AR{أَمْ} \AR{فِيدْ} \AR{أَبْ} & I am fed up \\
\AR{مَاكْ} \AR{آيْ} \AR{تَلْ} \AR{يُو} & Make I tell you & \AR{لِتْ} \AR{مِي} \AR{تِلْ} \AR{يُو} & Let me tell you \\

\AR{نَا} \AR{سُوْ} \AR{آيْ} \AR{سِي} \AR{آمْ} & Na so I see am & \AR{ذَتْسْ} \AR{هَوْ} \AR{آيْ} \AR{سَوْ} \AR{إِتْ} & That’s how I saw it \\
\AR{يُو} \AR{نَوْ} \AR{هِيِرْ} & You no hear & \AR{دِدْ} \AR{يُو} \AR{نُوتْ} \AR{هِيَرْ؟} & Did you not hear? \\
\AR{آيْ} \AR{دِي} \AR{فِينْدْ} & I dey find & \AR{آيْ} \AR{أَمْ} \AR{لُوكِنْجْ} \AR{فُورْ} & I am looking for \\
\AR{نَا} \AR{هُوْ} \AR{بِي} \AR{دَتْ؟} & Na who be that? & \AR{هُوْ} \AR{إِزْ} \AR{ذَتْ؟} & Who is that? \\
\AR{نَا} \AR{وِيتِنْ} \AR{بِي} \AR{دَتْ؟} & Na wetin be that? & \AR{وَتْ} \AR{إِزْ} \AR{ذَتْ؟} & What is that? \\
\AR{آيْ} \AR{نُو} \AR{سَابِي} & I no sabi & \AR{آيْ} \AR{دُونْتْ} \AR{أَنْدَرْسْتَانْدْ} & I don’t understand \\
\AR{نَا} \AR{سُوْ} \AR{إِهْ} \AR{بِي} & Na so e be & \AR{ذَتْسْ} \AR{هَوْ} \AR{إِتْ} \AR{إِزْ} & That’s how it is \\
\AR{آيْ} \AR{دِي} \AR{هِيِرْ} \AR{يُو} & I dey hear you & \AR{آيْ} \AR{أَنْدَرْسْتَانْدْ} \AR{يُو} & I understand you \\
\AR{يُو} \AR{فِيتْ} \AR{دُو} \AR{آمْ؟} & You fit do am? & \AR{كَانْ} \AR{يُو} \AR{دُو} \AR{إِتْ؟} & Can you do it? \\
\AR{آيْ} \AR{فِيتْ} \AR{تْرَاي} & I fit try & \AR{آيْ} \AR{كَانْ} \AR{تْرَاي} & I can try \\
\AR{نَا} \AR{سُوْ} \AR{لايْفْ} \AR{بِي} & Na so life be & \AR{ذَتْسْ} \AR{لايْفْ} & That’s life \\
\AR{آيْ} \AR{دِي} \AR{رَنْ} & I dey run & \AR{آيْ} \AR{أَمْ} \AR{إِنْ} \AR{أَ} \AR{هُورِي} & I am in a hurry \\
\AR{يُو} \AR{دُونْ} \AR{شُوبْ؟} & You don chop? & \AR{هَافْ} \AR{يُو} \AR{إِيتَنْ؟} & Have you eaten? \\
\AR{آيْ} \AR{نَوْ} \AR{وَاكَا} \AR{رِيتْ} & I no waka reach & \AR{آيْ} \AR{دِدْنْتْ} \AR{جِتْ} \AR{ذِيْرْ} & I didn’t get there \\
\AR{دِسْ} \AR{مَاتَرْ} \AR{هَارْدْ} & This matter hard & \AR{ذِسْ} \AR{إِزْ} \AR{دِفِيكَلْتْ} & This is difficult \\
\AR{نَا} \AR{سَمْتِينْ} \AR{دِي} \AR{رُونْ} & Na something dey wrong & \AR{سَمْثِينْجْ} \AR{إِزْ} \AR{رُونْجْ} & Something is wrong \\
\AR{آيْ} \AR{نَوْ} \AR{فِيتْ} \AR{كِلْ} \AR{مِيسِلْفْ} & I no fit kill myself & \AR{آيْ} \AR{كَانْتْ} \AR{أُوفَرْوُورْكْ} & I can’t overwork \\
\AR{يُو} \AR{دِي} \AR{يَانْ؟} & You dey yarn? & \AR{آرْ} \AR{يُو} \AR{تُوكِنْجْ؟} & Are you talking? \\
\AR{آيْ} \AR{دُونْ} \AR{يَانْ} \AR{آمْ} & I don yarn am & \AR{آيْ} \AR{أَلْرِيدِي} \AR{سِيدْ} \AR{إِتْ} & I already said it \\
\AR{مَاكْ} \AR{وِي} \AR{سِي} & Make we see & \AR{لِتْسْ} \AR{سِي} & Let’s see \\

\AR{يُو} \AR{نُو} \AR{بِيْتَرْ} & You know better & \AR{يُو} \AR{نُو} \AR{بِيْتَرْ} & You know better \\
\AR{آيْ} \AR{دِي} \AR{تِنْكْ} \AR{آمْ} & I dey think am & \AR{آيْ} \AR{أَمْ} \AR{ثِينْكِنْجْ} \AR{أَبَاوْتْ} \AR{إِتْ} & I am thinking about it \\
\AR{نَا} \AR{سَمْتِينْ} \AR{بِي} \AR{دِسْ} & Na something be this & \AR{ذِسْ} \AR{إِزْ} \AR{سَمْتِينْجْ} & This is something \\
\AR{يُو} \AR{نَوْ} \AR{دُو} \AR{آمْ} & You no do am & \AR{يُو} \AR{دِدْنْتْ} \AR{دُو} \AR{إِتْ} & You didn’t do it \\
\AR{آيْ} \AR{دُونْ} \AR{فُرْجِتْ} & I don forget & \AR{آيْ} \AR{هَافْ} \AR{فُرْجُوتِنْ} & I have forgotten \\
\AR{دِمْ} \AR{نَوْ} \AR{دِي} \AR{هِيِرْ} & Dem no dey hear & \AR{ذَيْ} \AR{أَرْ} \AR{نَاتْ} \AR{لِيسِنِنْجْ} & They are not listening \\
\AR{نَا} \AR{وِيتِنْ} \AR{آيْ} \AR{تَلْ} \AR{يُو} & Na wetin I tell you & \AR{ذَتْسْ} \AR{وَتْ} \AR{آيْ} \AR{تُولْدْ} \AR{يُو} & That’s what I told you \\
\AR{آيْ} \AR{دِي} \AR{كُولْ} & I dey cool & \AR{آيْ} \AR{أَمْ} \AR{كَالْمْ} & I am calm \\
\AR{يُو} \AR{دِي} \AR{فُورْسْ} \AR{آمْ} & You dey force am & \AR{يُو} \AR{أَرْ} \AR{فُورْسِنْجْ} \AR{إِتْ} & You are forcing it \\
\AR{نَا} \AR{نَاوْ} \AR{آيْ} \AR{نُو} & Na now I know & \AR{نَاوْ} \AR{آيْ} \AR{أَنْدَرْسْتَانْدْ} & Now I understand \\
\AR{آيْ} \AR{دِي} \AR{وِيتْ} & I dey wait & \AR{آيْ} \AR{أَمْ} \AR{وِيتِنْجْ} & I am waiting \\
\AR{يُو} \AR{نَوْ} \AR{جِتْ} \AR{سِنْسْ} & You no get sense & \AR{يُو} \AR{أَرْ} \AR{نَاتْ} \AR{ثِينْكِنْجْ} & You are not thinking \\
\AR{دِسْ} \AR{تِنْ} \AR{تَايَرْ} \AR{مِي} & This thing tire me & \AR{ذِسْ} \AR{مِيدْ} \AR{مِي} \AR{تَايَرْدْ} & This made me tired \\
\AR{نَا} \AR{مَايْ} \AR{فَالْتْ} & Na my fault & \AR{إِتْ} \AR{إِزْ} \AR{مَايْ} \AR{فَالْتْ} & It is my fault \\
\AR{آيْ} \AR{نَوْ} \AR{دِي} \AR{فِيلْ} \AR{آمْ} & I no dey feel am & \AR{آيْ} \AR{دُونْتْ} \AR{لَايْكْ} \AR{إِتْ} & I don’t like it \\
\AR{يُو} \AR{دِي} \AR{لُوكْ} \AR{مِي} & You dey look me & \AR{يُو} \AR{أَرْ} \AR{لُوكِنْجْ} \AR{أَتْ} \AR{مِي} & You are looking at me \\
\AR{نَا} \AR{سُوْ} \AR{آيْ} \AR{دِي} \AR{سِي} \AR{آمْ} & Na so I dey see am & \AR{ذَتْسْ} \AR{هَوْ} \AR{آيْ} \AR{سِي} \AR{إِتْ} & That’s how I see it \\
\AR{آيْ} \AR{دِي} \AR{تِيكْ} \AR{آمْ} & I dey take am & \AR{آيْ} \AR{وِلْ} \AR{دُو} \AR{إِتْ} & I will do it \\
\AR{دِمْ} \AR{دِي} \AR{كُولْ} \AR{مِي} & Dem dey call me & \AR{ذَيْ} \AR{أَرْ} \AR{كَالِنْجْ} \AR{مِي} & They are calling me \\
\AR{مَاكْ} \AR{وِي} \AR{رِيسْتْ} & Make we rest & \AR{لِتْسْ} \AR{رِيسْتْ} & Let’s rest \\
\AR{يُو} \AR{دِي} \AR{سْبِيكْ} \AR{بِيْجِنْ؟} & You dey speak pidgin? & \AR{دُو} \AR{يُو} \AR{سْبِيكْ} \AR{بِيْجِنْ؟} & Do you speak pidgin? \\
\AR{وِيتِنْ} \AR{لَنْجْوِيجْ} \AR{يُو} \AR{دِي} \AR{سْبِيكْ؟} & Wetin language you dey speak? & \AR{وَتْ} \AR{لَنْجْوِيجْ} \AR{دُو} \AR{يُو} \AR{سْبِيكْ؟} & What language do you speak? \\
\AR{يُو} \AR{دِي} \AR{سْبِيكْ} \AR{يَوْرُوبَا؟} & You dey speak Yoruba? & \AR{دُو} \AR{يُو} \AR{سْبِيكْ} \AR{يُورُوبَا؟} & Do you speak Yoruba? \\
\AR{يُو} \AR{دِي} \AR{سْبِيكْ} \AR{إِجْبُو؟} & You dey speak Igbo? & \AR{دُو} \AR{يُو} \AR{سْبِيكْ} \AR{إِجْبُو؟} & Do you speak Igbo? \\
\AR{يُو} \AR{دِي} \AR{سْبِيكْ} \AR{هَاوْسَا؟} & You dey speak Hausa? & \AR{دُو} \AR{يُو} \AR{سْبِيكْ} \AR{هَاوْسَا؟} & Do you speak Hausa? \\
\AR{يُو} \AR{إِنْجْلِيشْ} \AR{فَايْنْ} & Your English fine & \AR{يُو} \AR{إِنْجْلِيشْ} \AR{إِزْ} \AR{جْرِيتْ} & Your English is great \\
\AR{يُو} \AR{دِي} \AR{يَانْ} \AR{وِلْ} & You dey yarn well & \AR{يُو} \AR{سْبِيكْ} \AR{وِلْ} & You speak well \\
\AR{آيْ} \AR{لَايْكْ} \AR{هَاوْ} \AR{يُو} \AR{دِي} \AR{تَاكْ} & I like how you dey talk & \AR{آيْ} \AR{لَايْكْ} \AR{هَاوْ} \AR{يُو} \AR{سْبِيكْ} & I like how you speak \\
\AR{آيْ} \AR{نِيدْ} \AR{سُمْبُودِي} \AR{تُو} \AR{شُو} \AR{مِي} & I need somebody to show me & \AR{آيْ} \AR{نِيدْ} \AR{سَمْوَنْ} \AR{تُو} \AR{شُو} \AR{مِي} & I need someone to show me \\
\AR{شُو} \AR{مِي} \AR{هَاوْ} \AR{تُو} \AR{مَاكْ} \AR{فُوفُو} & Show me how to make fufu & \AR{شُو} \AR{مِي} \AR{هَاوْ} \AR{تُو} \AR{مَاكْ} \AR{فُوفُو} & Show me how to make fufu \\
\AR{شُو} \AR{مِي} \AR{هَاوْ} \AR{تُو} \AR{مَاكْ} \AR{إِجُوسِي} & Show me how to make egusi & \AR{شُو} \AR{مِي} \AR{هَاوْ} \AR{تُو} \AR{مَاكْ} \AR{إِجُوسِي} & Show me how to make egusi \\
\AR{آيْ} \AR{نَوْ} \AR{نُو} \AR{هَاوْ} \AR{تُو} \AR{كُوكْ} \AR{آمْ} & I no know how to cook am & \AR{آيْ} \AR{دُونْتْ} \AR{نُو} \AR{هَاوْ} \AR{تُو} \AR{كُوكْ} \AR{إِتْ} & I don’t know how to cook it \\
\AR{دِسْ} \AR{فُودْ} \AR{سْوِيتْ} \AR{وِلْ} \AR{وِلْ} & This food sweet well well & \AR{ذِسْ} \AR{فُودْ} \AR{إِزْ} \AR{فِيرِي} \AR{جُودْ} & This food is very good \\
\AR{آيْ} \AR{دِي} \AR{لِيرْنْ} & I dey learn & \AR{آيْ} \AR{أَمْ} \AR{لِيرْنِنْجْ} & I am learning \\
\AR{آيْ} \AR{دِي} \AR{تْرَاي} & I dey try & \AR{آيْ} \AR{أَمْ} \AR{تْرَايِنْجْ} & I am trying \\
\AR{تِيكْ} \AR{آمْ} \AR{سْمُولْ} \AR{سْمُولْ} & Take am small small & \AR{تِيكْ} \AR{إِتْ} \AR{سْلُولِي} & Take it slowly \\
\AR{نَا} \AR{سُوْ} \AR{دِمْ} \AR{تِيكْ} \AR{تُيچْ} \AR{مِي} & Na so dem take teach me & \AR{ذَتْسْ} \AR{هَوْ} \AR{ذَيْ} \AR{تُوتْ} \AR{مِي} & That’s how they taught me \\
\AR{آيْ} \AR{دِي} \AR{إِنْجُويْ} \AR{دِسْ} & I dey enjoy this & \AR{آيْ} \AR{أَمْ} \AR{إِنْجُويِنْجْ} & I am enjoying this \\
\AR{نَا} \AR{بِيْتَرْ} \AR{تِنْ} \AR{بِي} \AR{دِسْ} & Na better thing be this & \AR{ذِسْ} \AR{إِزْ} \AR{أَ} \AR{جُودْ} \AR{تِنْجْ} & This is a good thing \\
\AR{تَنْكْسْ} \AR{فُورْ} \AR{هِلْبْ} \AR{مِي} & Thanks for help me & \AR{ثَانْكْ} \AR{يُو} \AR{فُورْ} \AR{هِلْبِنْجْ} \AR{مِي} & Thank you for helping me \\

\end{longtable}

\section{Appendix M: Formal Specification of the Reverse Transliteration and English Normalization System}

\subsection{Overview}

The transliteration framework operates across three representational layers:

\[
\mathcal{A} \rightarrow \mathcal{E}_{\text{phon}} \rightarrow \mathcal{E}_{\text{std}}
\]

where:

\begin{itemize}
\item $\mathcal{A}$ denotes Arabic-script phonetic encoding,
\item $\mathcal{E}_{\text{phon}}$ denotes phonetic English transcription,
\item $\mathcal{E}_{\text{std}}$ denotes standardized orthographic English.
\end{itemize}

The system is therefore compositional:

\[
T = N \circ R
\]

where:

\[
R : \mathcal{A} \rightarrow \mathcal{E}_{\text{phon}}
\]

is the reverse transliteration function, and

\[
N : \mathcal{E}_{\text{phon}} \rightarrow \mathcal{E}_{\text{std}}
\]

is the normalization function.

\subsection{Stage One: Reverse Transliteration}

Let $\Sigma_A$ be the Arabic grapheme inventory used in the system.

Define a deterministic reading function:

\[
R : \Sigma_A^\ast \rightarrow \Sigma_{L}^\ast
\]

where $\Sigma_{L}$ denotes the Latin phonetic alphabet used for English approximation.

The function $R$ applies:

\begin{enumerate}
\item Grapheme-to-phoneme decoding,
\item Digraph resolution (e.g., \AR{ت}+\AR{ش} $\mapsto$ \textit{ch}),
\item Vowel interpretation under marked diacritics,
\item Sequential concatenation.
\end{enumerate}

Under full vowel marking and digraph recognition, $R$ is operationally deterministic.

\subsection{Stage Two: Orthographic Normalization}

The phonetic English layer does not guarantee standard spelling. Therefore define:

\[
N : \Sigma_{L}^\ast \rightarrow \Sigma_{E}^\ast
\]

where $\Sigma_{E}$ is the standardized English orthographic system.

The normalization operator performs:

\begin{enumerate}
\item Lexical correction (e.g., \textit{utopianizm} $\rightarrow$ \textit{utopianism}),
\item Morphological restoration (e.g., possessives, suffix normalization),
\item Conventional capitalization,
\item Punctuation smoothing.
\end{enumerate}

The operator $N$ is not purely phonemic; it references an English lexicon and morphological rules.

\subsection{Compositional Structure}

The full recovery mapping is:

\[
T = N(R(\cdot))
\]

Thus:

\[
\mathcal{A} \xrightarrow{R} \mathcal{E}_{\text{phon}} \xrightarrow{N} \mathcal{E}_{\text{std}}
\]

The first mapping is phonological; the second is orthographic and lexical.

\subsection{Determinism and Ambiguity}

The mapping $R$ is deterministic under the following conditions:

\[
\text{Full vowel marking} \land \text{Digraph priority parsing}
\]

The mapping $N$ may admit multiple outputs if phonetic spellings correspond to multiple lexical candidates. In practice, contextual constraints restrict ambiguity.

\subsection{Loss and Recovery}

Let:

\[
L = R^{-1}
\]

be forward encoding into Arabic script.

Then:

\[
R(L(x)) = x
\]

holds when $x$ is phonemically spelled.

However:

\[
N(R(L(x))) = x
\]

holds only if $x$ conforms to normalized English spelling conventions.

Thus the system preserves phonetic structure strictly and orthographic convention conditionally.

\subsection{Structural Interpretation}

The transliteration architecture is layered:

\begin{itemize}
\item Layer 1: Acoustic encoding (Arabic graphemes),
\item Layer 2: Phonetic English reconstruction,
\item Layer 3: Standardized literary English.
\end{itemize}

This confirms that the Arabic-script text encodes sound rather than etymology. Conventional English spelling is a secondary reconstruction layer, not the primary representation.

\subsection{Conclusion}

The system is:

\[
\text{Phonologically grounded}
\]

\[
\text{Algorithmically compositional}
\]

\[
\text{Reversible at the phonetic level}
\]

\[
\text{Lexically normalizable at the literary level}
\]

The separation of $R$ and $N$ ensures conceptual clarity between acoustic transcription and orthographic standardization.

\section{Appendix N: Finite-State Formalization of the Transliteration System}

\subsection{Overview}

The transliteration architecture may be modeled as a cascade of deterministic finite-state transducers (FSTs):

\[
\mathcal{A} \xrightarrow{R} \mathcal{E}_{\text{phon}} \xrightarrow{N} \mathcal{E}_{\text{std}}
\]

Each mapping is computable in linear time with respect to input length.

\subsection{Alphabet Definitions}

Let:

\[
\Sigma_A = \{\text{Arabic graphemes used in the system}\}
\]

\[
\Sigma_P = \{\text{phoneme symbols for English/Pidgin}\}
\]

\[
\Sigma_E = \{\text{standard English orthographic symbols}\}
\]

All alphabets are finite.

\subsection{Transducer 1: Reverse Transliteration}

Define transducer:

\[
R = (Q_R, \Sigma_A, \Sigma_P, \delta_R, q_0, F_R)
\]

where:

\begin{itemize}
\item $Q_R$ is a finite set of states,
\item $\delta_R$ is the transition function,
\item $q_0$ is the start state,
\item $F_R$ is the set of accepting states.
\end{itemize}

The transition function implements:

\begin{enumerate}
\item Single-letter grapheme-to-phoneme mapping,
\item Priority digraph detection (e.g., \AR{ت}+\AR{ش} $\rightarrow$ /tʃ/),
\item Vowel diacritic resolution,
\item Sequential concatenation.
\end{enumerate}

Digraph priority is implemented via a two-state lookahead mechanism:

\[
\delta_R(q_0, \AR{ت}) = q_{T}
\]

\[
\delta_R(q_{T}, \AR{ش}) = q_0 \; \text{with output } /tʃ/
\]

Otherwise:

\[
\delta_R(q_{T}, x \neq \AR{ش}) = q_0 \; \text{with output } /t/ \text{ and reprocess } x
\]

This ensures maximal matching without nondeterminism.

\subsection{Transducer 2: Normalization Operator}

Define lexical normalization transducer:

\[
N = (Q_N, \Sigma_P, \Sigma_E, \delta_N, q'_0, F_N)
\]

This operator performs:

\begin{itemize}
\item Phonetic-to-orthographic replacement,
\item Morphological reconstruction,
\item Capitalization rules,
\item Punctuation smoothing.
\end{itemize}

Unlike $R$, $N$ references a finite lexicon:

\[
\mathcal{L} \subset \Sigma_P^\ast
\]

The normalization rule may be written as:

\[
\delta_N(q, w_{\text{phon}}) = (q', w_{\text{std}})
\]

whenever:

\[
w_{\text{phon}} \in \mathcal{L}
\]

If no lexical rule applies, identity mapping is used.

\subsection{Composed System}

The composed transliteration machine is:

\[
T = N \circ R
\]

Since finite-state transducers are closed under composition:

\[
T
\]

is also a finite-state transducer.

Thus the entire Arabic-to-English system is regular and computable.

\subsection{Time Complexity}

For input string length $n$:

\[
R = O(n)
\]

\[
N = O(n)
\]

Therefore:

\[
T = O(n)
\]

The system is linear in input length.

\subsection{Reversibility Conditions}

Let $L$ be the forward encoding operator:

\[
L : \Sigma_P^\ast \rightarrow \Sigma_A^\ast
\]

Then:

\[
R(L(x)) = x
\]

for fully vowel-marked phonemic input.

However:

\[
N(R(L(x))) = x
\]

holds only when $x$ is lexically normalized.

Thus phonemic reversibility is strict; orthographic reversibility is conditional.

\subsection{Structural Significance}

This formalization confirms that:

\begin{itemize}
\item The transliteration system is algorithmically definable.
\item No context-free machinery is required.
\item Digraph handling does not introduce nondeterminism.
\item The architecture is compositional and modular.
\end{itemize}

\subsection{Conclusion}

The Arabic-script Pidgin and English transliteration framework constitutes a cascade of deterministic finite-state transducers. It is therefore:

\[
\text{Regular}
\]

\[
\text{Computationally tractable}
\]

\[
\text{Formally compositional}
\]

This places the system within the standard computational paradigm used in modern orthographic engineering and digital linguistic modeling.

\newpage
\newpage
\begin{thebibliography}{99}

\bibitem{AbdulazizOsinde1997}
Abdulaziz, Mohamed H., and Ken Osinde (1997).
Sheng and Engsh: Development of Mixed Codes among the Urban Youth in Kenya.
\textit{International Journal of the Sociology of Language} 125, 43--63.

\bibitem{Adegbija2004}
Adegbija, Efurosibina (2004).
The Domestication of English in Nigeria.
In: Edgar W. Schneider et al. (eds.),
\textit{A Handbook of Varieties of English, Volume 1}.
Berlin: Mouton de Gruyter, 20--43.

\bibitem{Aikhenvald2006}
Aikhenvald, Alexandra Y. (2006).
\textit{Serial Verb Constructions in Typological Perspective}.
Oxford: Oxford University Press.

\bibitem{AikhenvaldDixon2006}
Aikhenvald, Alexandra Y., and R. M. W. Dixon (eds.) (2006).
\textit{Serial Verb Constructions: A Cross-Linguistic Typology}.
Oxford: Oxford University Press.

\bibitem{Akinnaso1995}
Akinnaso, F. Niyi (1995).
Literacy, Language and Education in Nigeria.
\textit{Comparative Education Review} 39(1), 68--92.

\bibitem{Ayorinde2013}
Ayorinde, Christine (2013).
Language Policy and the Development of Nigerian Pidgin.
In: Ore Yusuf (ed.),
\textit{Language and Society in Nigeria}.
Ibadan: University Press, 201--220.

\bibitem{Bakker2014}
Bakker, Peter (2014).
Pidgins.
In: Silvia Luraghi and Claudia Parodi (eds.),
\textit{The Bloomsbury Companion to Syntax}.
London: Bloomsbury, 109--130.

\bibitem{BaldiCuly2005}
Baldi, Philip, and Christopher Culy (eds.) (2005).
\textit{Computational Approaches to African Languages}.
Stanford: CSLI Publications.

\bibitem{BeesleyKarttunen2003}
Beesley, Kenneth R., and Lauri Karttunen (2003).
\textit{Finite State Morphology}.
Stanford: CSLI Publications.

\bibitem{Besacier2014}
Besacier, Laurent, et al. (2014).
Automatic Speech Recognition for Under-Resourced African Languages.
\textit{Speech Communication} 56, 85--100.

\bibitem{Bickerton1981}
Bickerton, Derek (1981).
\textit{Roots of Language}.
Ann Arbor: Karoma.

\bibitem{Bickerton1984}
Bickerton, Derek (1984).
The Language Bioprogram Hypothesis.
\textit{Behavioral and Brain Sciences} 7(2), 173--221.

\bibitem{BirdSimons2003}
Bird, Steven, and Gary Simons (2003).
Seven Dimensions of Portability for Language Documentation and Description.
\textit{Language} 79(3), 557--582.

\bibitem{BlanchonLafkioui2017}
Blanchon, Jean A., and Mena Lafkioui (eds.) (2017).
\textit{African Arabic and Contact Linguistics}.
Leiden: Brill.

\bibitem{Blench2006}
Blench, Roger (2006).
\textit{The Languages of Africa}.
Cambridge: Cambridge University Press.

\bibitem{Blench2011}
Blench, Roger (2011).
The Status of Nigerian Pidgin.
\textit{Cambridge Working Papers in Linguistics} 3, 1--27.

\bibitem{Bondarev2017}
Bondarev, Dmitry (2017).
Ajami in West Africa.
In: Fallou Ngom (ed.),
\textit{Muslims Beyond the Arab World: The Odyssey of \textit{Ajam\={i}} and the Mur\={i}diyya}.
Oxford: Oxford University Press, 27--53.

\bibitem{Bybee2010}
Bybee, Joan (2010).
\textit{Language, Usage and Cognition}.
Cambridge: Cambridge University Press.

\bibitem{Cisse2018}
Cissé, Mamadou (2018).
Ajami Writing Systems in West Africa.
\textit{Islamic Africa} 9(2), 135--160.

\bibitem{Comrie1985}
Comrie, Bernard (1985).
\textit{Tense}.
Cambridge: Cambridge University Press.

\bibitem{Coulmas2003}
Coulmas, Florian (2003).
\textit{Writing Systems: An Introduction to Their Linguistic Analysis}.
Cambridge: Cambridge University Press.

\bibitem{Daniels2017}
Daniels, Peter T. (2017).
Script Reform and Orthographic Standardization in Africa.
In: Badreddine Arfi (ed.),
\textit{Writing Systems and Orthography in Africa}.
Dakar: CODESRIA, 45--78.

\bibitem{DanielsBright1996}
Daniels, Peter T., and William Bright (eds.) (1996).
\textit{The World's Writing Systems}.
Oxford: Oxford University Press.

\bibitem{Dwyer1998}
Dwyer, David (1998).
The Use of Arabic Script for African Languages.
In: Karin Barber (ed.),
\textit{Africa's Hidden Histories}.
Bloomington: Indiana University Press, 231--252.

\bibitem{Faraclas1996}
Faraclas, Nicholas (1996).
\textit{Nigerian Pidgin}.
London: Routledge.

\bibitem{Faraclas2013}
Faraclas, Nicholas (2013).
Nigerian Pidgin.
In: Susanne Maria Michaelis et al. (eds.),
\textit{The Survey of Pidgin and Creole Languages}.
Oxford: Oxford University Press, 435--445.

\bibitem{Furniss2004}
Furniss, Graham (2004).
\textit{Literature in Hausa}.
Edinburgh: Edinburgh University Press.

\bibitem{Gambari2015}
Gambari, A. I. (2015).
Ajami Manuscripts and Literacy in Northern Nigeria.
\textit{Sudanic Africa} 26, 87--110.

\bibitem{Hasselblatt2006}
Hasselblatt, Cornelius (2006).
Orthography as a Linguistic Discipline.
\textit{Written Language and Literacy} 9(1), 1--23.

\bibitem{Hasselbring2007}
Hasselbring, Sue (2007).
Hausa Ajami Manuscripts and Orthographic Variation.
\textit{Sudanic Africa} 18, 105--132.

\bibitem{Hockett1958}
Hockett, Charles F. (1958).
\textit{A Course in Modern Linguistics}.
New York: Macmillan.

\bibitem{Holm1988}
Holm, John (1988).
\textit{Pidgins and Creoles, Volume I: Theory and Structure}.
Cambridge: Cambridge University Press.

\bibitem{Holm1989}
Holm, John (1989).
\textit{Pidgins and Creoles, Volume II: Reference Survey}.
Cambridge: Cambridge University Press.

\bibitem{Huber1999}
Huber, Magnus (1999).
\textit{Ghanaian Pidgin English in Its West African Context}.
Amsterdam: John Benjamins.

\bibitem{Huber2008}
Huber, Magnus, and Viveka Velupillai (eds.) (2008).
\textit{Synchronic and Diachronic Perspectives on Contact Languages}.
Amsterdam: John Benjamins.

\bibitem{Hunwick1999}
Hunwick, John (1999).
\textit{Arabic Literature of Africa, Volume 4}.
Leiden: Brill.

\bibitem{Hyman2003}
Hyman, Larry M. (2003).
Segmental Phonology.
In: Derek Nurse and Gérard Philippson (eds.),
\textit{The Bantu Languages}.
London: Routledge, 42--58.

\bibitem{Hyman2007}
Hyman, Larry M. (2007).
Universals of Tone Rules: 30 Years Later.
In: Tomas Riad and Carlos Gussenhoven (eds.),
\textit{Tones and Tunes, Volume 1}.
Berlin: Mouton de Gruyter, 1--34.

\bibitem{KaplanKay1994}
Kaplan, Ronald M., and Martin Kay (1994).
Regular Models of Phonological Rule Systems.
\textit{Computational Linguistics} 20(3), 331--378.

\bibitem{Karan2014}
Karan, Elke (2014).
Orthography Development and Reform.
In: James W. Tollefson and Miguel Pérez-Milans (eds.),
\textit{The Oxford Handbook of Language Policy and Planning}.
Oxford: Oxford University Press.

\bibitem{Karanetal2010}
Karan, Elke, Susan Malone, and John Smalley (2010).
\textit{Developing Orthographies for Unwritten Languages}.
Dallas: SIL International.

\bibitem{Kouwenberg2008}
Kouwenberg, Silvia, and John Victor Singler (eds.) (2008).
\textit{The Handbook of Pidgin and Creole Studies}.
Malden, MA: Wiley-Blackwell.

\bibitem{Labov1972}
Labov, William (1972).
\textit{Sociolinguistic Patterns}.
Philadelphia: University of Pennsylvania Press.

\bibitem{LewisSimonsFennig2016}
Lewis, M. Paul, Gary F. Simons, and Charles D. Fennig (eds.) (2016).
\textit{Ethnologue: Languages of the World}, 19th edn.
Dallas: SIL International.

\bibitem{Lupke2011}
Lüpke, Friederike (2011).
Orthography Development.
In: Peter Austin and Julia Sallabank (eds.),
\textit{The Cambridge Handbook of Endangered Languages}.
Cambridge: Cambridge University Press, 312--336.

\bibitem{McWhorter1998}
McWhorter, John (1998).
Identifying the Creole Prototype.
\textit{Language} 74(4), 788--818.

\bibitem{McWhorter2005}
McWhorter, John (2005).
\textit{Defining Creole}.
Oxford: Oxford University Press.

\bibitem{Michaelis2013}
Michaelis, Susanne Maria, et al. (eds.) (2013).
\textit{The Survey of Pidgin and Creole Languages}, 3 vols.
Oxford: Oxford University Press.

\bibitem{Mufwene2001}
Mufwene, Salikoko S. (2001).
\textit{The Ecology of Language Evolution}.
Cambridge: Cambridge University Press.

\bibitem{Mufwene2008}
Mufwene, Salikoko S. (2008).
\textit{Language Evolution: Contact, Competition, and Change}.
London: Continuum.

\bibitem{Ngom2010}
Ngom, Fallou (2010).
Ajami Scripts in the Senegalese Speech Community.
\textit{Journal of Arabic and Islamic Studies} 10, 1--23.

\bibitem{Ngom2016}
Ngom, Fallou (2016).
\textit{Muslims Beyond the Arab World: The Odyssey of \textit{Ajam\={i}} and the Mur\={i}diyya}.
Oxford: Oxford University Press.

\bibitem{Ngugi1986}
Ngũgĩ wa Thiong'o (1986).
\textit{Decolonising the Mind}.
London: James Currey.

\bibitem{Oyetade2006}
Oyetade, S. O. (2006).
Attitudes to Foreign Languages and Indigenous Languages in Nigeria.
\textit{Sociolinguistic Studies} 1(1), 19--38.

\bibitem{Oyetade2008}
Oyetade, Solomon (2008).
Language Planning in West Africa.
In: Andrew Simpson (ed.),
\textit{Language and National Identity in Africa}.
Oxford: Oxford University Press, 165--187.

\bibitem{Philippson2012}
Philippson, Gérard (2012).
Writing African Languages.
In: Bernd Heine and Derek Nurse (eds.),
\textit{The Oxford Handbook of African Languages}.
Oxford: Oxford University Press, 781--799.

\bibitem{Robinson2006}
Robinson, Clinton (2006).
Language Policy and Literacy in Africa.
\textit{International Journal of Educational Development} 26(4), 379--391.

\bibitem{Rogers2005}
Rogers, Henry (2005).
\textit{Writing Systems: A Linguistic Approach}.
Oxford: Blackwell.

\bibitem{Sanneh1989}
Sanneh, Lamin (1989).
\textit{Translating the Message}.
Maryknoll: Orbis Books.

\bibitem{Schmidt2008}
Schmidt, Annette (2008).
Computational Tools for African Orthographies.
In: Silvia Kouwenberg and John Victor Singler (eds.),
\textit{The Handbook of Pidgin and Creole Studies}.
Malden, MA: Wiley-Blackwell, 570--589.

\bibitem{Schneider2007}
Schneider, Edgar W. (2007).
\textit{Postcolonial English: Varieties Around the World}.
Cambridge: Cambridge University Press.

\bibitem{SIL2004}
SIL International (2004).
\textit{Orthography Development Manual}.
Dallas: SIL International.

\bibitem{Smalley1964}
Smalley, William A. (1964).
Orthography Studies.
\textit{Anthropological Linguistics} 6(2), 1--11.

\bibitem{Smalley1988}
Smalley, William A. (1988).
Orthography as a Social Practice.
\textit{Language in Society} 17(3), 377--393.

\bibitem{Suleiman2013}
Suleiman, Yasir (2013).
\textit{Arabic in the Fray: Language Ideology and Cultural Politics}.
Edinburgh: Edinburgh University Press.

\bibitem{ThomasonKaufman1988}
Thomason, Sarah G., and Terrence Kaufman (1988).
\textit{Language Contact, Creolization, and Genetic Linguistics}.
Berkeley: University of California Press.

\bibitem{Tosco2000}
Tosco, Mauro (2000).
Is Nigerian Pidgin English a Creole?
\textit{Journal of Pidgin and Creole Languages} 15(1), 1--25.

\bibitem{Vydrin2012}
Vydrin, Valentin (2012).
Scripts and Orthographies in West Africa.
\textit{Mandenkan} 48, 3--30.

\bibitem{Winford2003}
Winford, Donald (2003).
\textit{An Introduction to Contact Linguistics}.
Malden, MA: Blackwell.

\bibitem{Winford2008}
Winford, Donald (2008).
Creoles and Contact.
In: Silvia Kouwenberg and John Victor Singler (eds.),
\textit{The Handbook of Pidgin and Creole Studies}.
Malden, MA: Wiley-Blackwell, 3--41.

\bibitem{Zima1974}
Zima, Petr (1974).
\textit{Pidginization and Creolization: The Case of West African Pidgin English}.
Munich: Wilhelm Fink Verlag.

\bibitem{Zwartjes2011}
Zwartjes, Otto (2011).
\textit{Portuguese Missionary Grammars in Asia, Africa and Brazil}.
Amsterdam: John Benjamins.

\bibitem{Affia2023}
Affia, Uchenna (2023).
Nigerian Pidgin: The Identity of a Nigerian Away from Home.
Unpublished thesis, York University, Toronto.

\bibitem{AjaKilani2024}
Ajayi, Taiwo, and Mukhtar Kilani (2024).
Making Nigerian Pidgin English an Official Language in Nigeria: Prospects and Challenges.
\textit{Journal of Language and Linguistic Studies} 20(2), 1--22.

\bibitem{AmadiKilani2021}
Amadi, Chukwuemeka, and Mukhtar Kilani (2021).
Nigerian Pidgin English: An Effective Communicative Device for Vocabulary Development.
\textit{Sapientia Global Journal of Arts, Humanities and Development Studies} 4(3), 89--102.

\bibitem{Ehondor2020}
Ehondor, Beryl Annette (2020).
Nigerian Pidgin English: A Cultural Universal for National Communication and Policy Enactment.
\textit{Pan-Atlantic University Working Papers in Communication and Media Studies} 1(1), 1--18.

\bibitem{Frackiewicz2019}
Fr\k{a}ckiewicz, Marta (2019).
Nigerian Pidgin English Phraseology in the Context of Areal Influences.
\textit{Studies in African Languages and Cultures} 53, 5--34.

\bibitem{OjoKilani2020}
Ojo, Emmanuel, and Mukhtar Kilani (2020).
The Nigerian Pidgin English: A Tool for National Integration.
\textit{Journal of the Linguistic Association of Nigeria} 23(1), 44--59.

\end{thebibliography}

\end{document}
